% main.tex -- Paper 3: energy and cost per query of CPU and GPU dataframe engines.
%
% Target venue: ACM e-Energy 2027, full paper: at most 10 pages in 9-pt ACM sigconf, excluding
% references and appendices; double-blind with non-destructive anonymization at submission.
% e-Energy 2027 CFP (energy.acm.org/conferences/eenergy/2027/pages/cfp.php; read by venue-scout through a
% browser on 2026-11-23): full papers "up to 10 pages in 9-point ACM double-column format, excluding
% references and appendices"; camera-ready gets one additional page; the appendix "is not constrained;
% however, it may only be used for the purpose of justifying the technical correctness of their claims
% stated in the body of the paper"; double-blind with "only non-destructive anonymization"; own past
% work cited "in the third person"; "Remove acknowledgment of identifying names and funding sources."
% The CFP says nothing about code, data or pre-registration links (withheld here in the anonymized build).
% CCS concepts: verified 2026-11-23 (see the CCSXML block). Keywords: the author's choice (ACM requires them).
%
% Builds: this file is the ANONYMIZED submission build (default). camera-ready.tex defines
% \pthreecameraready and inputs this file: author, acknowledgments, funding and credits, and
% the self-identifying links (repository, Zenodo DOIs, commits) appear only there
% (macros.tex: \camera, \selfdoi, \repourl, \lockdoi, \lockcommit, \citelock).
%
% STATUS: skeleton written 2026-11-23 while the campaign ran, before any campaign value was seen; results
% filled in 2026-11-24 from the frozen analysis output (results/paper3-analysis/stats.json, analysis code
% sha256 80ac76fb...). Methods follow docs/PAPER3_ENERGY_DESIGN.md (locked 2026-11-16; amendments A1--A11).
% Every campaign result is a \stat-family macro whose value comes from stats.json via `make stats`; the
% judgments are the analysis's own class words. Pre-lock probe values (synthetic load; design 2.1-2.4, 4)
% are method validation and are typed.
%
% Amendment A5 (GH200 and B200 not run; GCP funding; start order) was checked against the committed
% text (commit 39b8f7b) on 2026-11-23.

\ifdefined\pthreecameraready
  \documentclass[sigconf]{acmart}
\else
  \documentclass[sigconf,anonymous,review]{acmart}
\fi

\usepackage{multirow}
\usepackage{enumitem}
\usepackage{xspace}
\usepackage{seqsplit}  % breakable TODO keys in draft builds
\usepackage{array}
\newcolumntype{R}[1]{>{\raggedright\arraybackslash}p{#1}}  % ragged-right paragraph column
% macros.tex -- Paper 3: statistics lookup, anonymization switch, figure, notice and notation macros.
%
% Numbers from the campaign are never typed in. Every result is written as \stat{key} (running
% text) or \tstat{key} (table cells), where key is a dotted stats.json key such as
% x1.h100.baseline.sf300.energy_total.estimate. The \stat machinery is Paper 1's
% (../paper/macros.tex): a stats file generated/stats.tex with one \defstat{key}{value} per key
% and \statsfinaltrue. In final mode a missing key is a hard LaTeX error; in draft mode (no
% generated/stats.tex, as now) the key is printed in a yellow TODO box, and table cells print a
% red TODO (the cell's key pattern is in the table's caption comment).
%
% TODO(results): the key names here are provisional. They must match the stats.json written by
% the Paper 3 analysis code (not in the repository on 2026-11-23); a generator like
% ../paper/tools/stats_to_tex.py must produce generated/stats.tex from it.
%
% Pre-lock method-validation numbers (probes P1-P7, G8 probe; design sections 2.1-2.4, 4) are
% typed directly: they are synthetic-load QA values fixed before the lock, not campaign outcomes.

\makeatletter
\newif\ifstatsfinal
\statsfinalfalse
\newcommand{\defstat}[2]{\expandafter\def\csname stat@#1\endcsname{#2}}
% Breakable (seqsplit) so that long keys do not overfill two-column lines in draft builds.
\newcommand{\statplaceholder}[1]{%
  {\color{red}\tiny\ttfamily[TODO:\expandafter\seqsplit\expandafter{\detokenize{#1}}]}}
\newcommand{\stat@missing}[1]{%
  \PackageError{stats}{Missing statistic '#1' in generated/stats.tex}%
    {Re-run the stats generator; the analysis output lacks this key.}}
\newcommand{\stat}[1]{%
  \ifcsname stat@#1\endcsname
    \csname stat@#1\endcsname
  \else
    \ifstatsfinal\stat@missing{#1}\else\statplaceholder{#1}\fi
  \fi}
% Compact form for table cells.
\newcommand{\tstat}[1]{%
  \ifcsname stat@#1\endcsname
    \csname stat@#1\endcsname
  \else
    \ifstatsfinal\stat@missing{#1}\else\textcolor{red}{\tiny TODO}\fi
  \fi}
% "estimate [ci_low, ci_high]" for a ratio prefix such as x1.h100.baseline.sf300.energy_total
\newcommand{\ratioci}[1]{%
  \stat{#1.estimate}\,{\footnotesize[\stat{#1.ci_low},\,\stat{#1.ci_high}]}}
% "estimate x (95% CI ci_low--ci_high)" for running text
\newcommand{\ratiotext}[1]{%
  \stat{#1.estimate}$\times$ (95\%~CI \stat{#1.ci_low}--\stat{#1.ci_high})}
% Table form of \ratioci: one red TODO in draft mode instead of three.
\newcommand{\tratioci}[1]{%
  \ifcsname stat@#1.estimate\endcsname
    \tstat{#1.estimate}\,{\scriptsize[\tstat{#1.ci_low},\,\tstat{#1.ci_high}]}%
  \else
    \ifstatsfinal\stat@missing{#1.estimate}\else\textcolor{red}{\tiny TODO}\fi
  \fi}
\makeatother

% Two value files (make stats; make record):
%  generated/stats.tex  -- copied unchanged from analyze_paper3.py's OUT/stats.tex (one \defstat per
%                          scalar of its flattened stats.json; \syntheticstatstrue on a synthetic tree).
%                          It does not set \statsfinaltrue, so its presence switches final mode on here.
%  generated/record.tex -- hand-kept \defstat lines for values from committed records, not campaign
%                          outcomes: amendment A6 (date, commit, sha256) and launch-record facts
%                          (campaign.<host>.driver, .cpu_model, .first_session_utc, ...).
\newif\ifsyntheticstats
\IfFileExists{generated/record.tex}{\input{generated/record.tex}}{}
\IfFileExists{generated/stats.tex}{\input{generated/stats.tex}\statsfinaltrue}{}

% Watermarks: DRAFT while no stats file is present; SYNTHETIC when the stats come from a synthetic
% tree. No draft or synthetic page can be mistaken for a results page.
\newcommand{\watermarktext}{%
  \ifsyntheticstats SYNTHETIC -- NOT RESULTS\else\ifstatsfinal\else DRAFT -- NO RESULTS\fi\fi}
\AddToHook{shipout/background}{%
  \ifsyntheticstats\def\pthreewmon{1}\else\ifstatsfinal\def\pthreewmon{0}\else\def\pthreewmon{1}\fi\fi
  \ifnum\pthreewmon=1
  \put(0.5\paperwidth,-0.5\paperheight){\makebox(0,0){\rotatebox{45}{%
    \scalebox{4}{\sffamily\bfseries\textcolor{red!15}{\watermarktext}}}}}\fi}

% ---------------------------------------------------------------- anonymization switch
% main.tex is the anonymized submission build (the default). camera-ready.tex defines
% \pthreecameraready and inputs main.tex: that build shows the author, the acknowledgments,
% the funding and credits, and every self-identifying link (repository, Zenodo DOIs, commits).
% Self-citations of Papers 1 and 2 stay in the third person in both builds (non-destructive
% anonymization); only their DOIs are withheld.
\newif\ifanon
\ifdefined\pthreecameraready\anonfalse\else\anontrue\fi
% \camera{camera-ready text}{anonymized text}
\newcommand{\camera}[2]{\ifanon#2\else#1\fi}
\newcommand{\anonmark}{[withheld for review]}
\newcommand{\selfdoi}[1]{\camera{\href{https://doi.org/#1}{doi:#1}}{DOI withheld for review}}
\newcommand{\repourl}{\camera{\url{https://github.com/ejiod2025/gpu-spark-polars-study}}{\anonmark}}
\newcommand{\lockdoi}{\camera{\href{https://doi.org/10.5281/zenodo.24284712}{doi:10.5281/zenodo.24284712}}{\anonmark}}
\newcommand{\lockcommit}{\camera{\texttt{bd42401}}{\anonmark}}
\newcommand{\citelock}{\camera{\cite{ejiofor2026p3lock}}{\cite{p3lockanon}}}

% Author-supplied detail still missing, and results-not-yet-available markers for draft prose
% (each also flagged with a "% TODO(...)" comment in the source).
\newcommand{\todoplaceholder}[1]{\textcolor{red}{[TODO: #1]}}
\newcommand{\todoverify}[1]{\textcolor{red}{[TODO-VERIFY: #1]}}
\newcommand{\pending}[1]{\textcolor{gray}{[TODO(results): #1]}}

% Figures are copied into figures/ once the analysis exists. Until then a labelled placeholder
% box keeps the layout and cross-references working. #3 is the planned height.
\newcommand{\studyfig}[3]{%
  \IfFileExists{figures/#1.pdf}{\includegraphics[width=#2]{figures/#1.pdf}}{%
  \fbox{\parbox[c][#3][c]{0.92\linewidth}{\centering\footnotesize\ttfamily
    TODO(results): figure pending\\\detokenize{#1}}}}}

% TPC fair-use disclaimer, Paper 1's wording (../paper/macros.tex \tpcdisclaimer) with the author's
% workload name for Paper 3 (design sec. 8, 2026-11-11) in place of Paper 1's \workloadname.
\newcommand{\workloadname}{TPC-H-derived workload}
\newcommand{\tpcdisclaimer}{The \workloadname{} is derived from the TPC-H
  Benchmark and as such is not comparable to published TPC-H Benchmark results,
  as the \workloadname{} results do not comply with the TPC-H Benchmark
  Specification.}

% Engine names (Paper 1) and hosts
\newcommand{\eng}[1]{\texttt{#1}}
\newcommand{\polarscpu}{Polars CPU}
\newcommand{\polarsgpu}{Polars GPU}
\newcommand{\sparkcpu}{Spark CPU}
\newcommand{\sparkgpu}{Spark GPU}
\newcommand{\q}[1]{\texttt{#1}}
\newcommand{\host}[1]{\textsf{#1}}
\newcommand{\GA}{\host{G-A100}}
\newcommand{\GH}{\host{G-H100}}
\newcommand{\CG}{\host{C-GCP}}
\newcommand{\CV}{\host{C-VM}}

% Notation
\newcommand{\scope}[1]{\textsf{#1}}          % scope tags: S1-counter, S2-pinned, ...
\newcommand{\tmin}{t_{\min}}
\newcommand{\Pidle}{P_{\mathrm{idle}}}
\newcommand{\Eabove}{E_{\mathrm{above}}}
\newcommand{\Dstar}{\Delta^{*}}
\newcommand{\NOne}{\textsf{N1}}

% Compact contrast cell for wide tables: "estimate^class". The class is the analysis's word
% (design 5.3; stats.tex escapes "_" as "\_"), shown as G (gpu_better), C (cpu_better),
% = (no_difference), ? (inconclusive) or -- (not_estimable); one red TODO in draft mode.
\makeatletter
\edef\xcls@gpu{\detokenize{gpu\_better}}
\edef\xcls@cpu{\detokenize{cpu\_better}}
\edef\xcls@nod{\detokenize{no\_difference}}
\edef\xcls@inc{\detokenize{inconclusive}}
\newcommand{\xclass}[1]{%
  \ifcsname stat@#1\endcsname
    \edef\xcls@v{\expandafter\expandafter\expandafter\detokenize\expandafter\expandafter
      \expandafter{\csname stat@#1\endcsname}}%
    \ifx\xcls@v\xcls@gpu G\else\ifx\xcls@v\xcls@cpu C\else
    \ifx\xcls@v\xcls@nod =\else\ifx\xcls@v\xcls@inc ?\else --\fi\fi\fi\fi
  \else
    \ifstatsfinal\stat@missing{#1}\else\textcolor{red}{\tiny TODO}\fi
  \fi}
\newcommand{\xcell}[1]{%
  \ifcsname stat@#1.estimate\endcsname
    \tstat{#1.estimate}\textsuperscript{\xclass{#1.class}}%
  \else
    \ifstatsfinal\stat@missing{#1.estimate}\else\textcolor{red}{\tiny TODO}\fi
  \fi}
\makeatother


\setlist{leftmargin=*,itemsep=1pt,topsep=2pt}
\setlength{\emergencystretch}{3em}

% ---------------------------------------------------------------- ACM metadata
% Conference: e-Energy 2027 CFP (https://energy.acm.org/conferences/eenergy/2027/pages/cfp.php, read by
% venue-scout through a browser on 2026-11-23). TODO-VERIFY(venue): the rights block, set at acceptance.
\setcopyright{acmlicensed}
\copyrightyear{2027}
\acmYear{2027}
\acmDOI{XXXXXXX.XXXXXXX}
\acmConference[e-Energy '27]{The 18th ACM International Conference on Future and Sustainable
  Energy Systems}{June 5--11, 2027}{Atlanta, GA, USA}
\acmISBN{978-X-XXXX-XXXX-X/2027/XX}

\begin{document}

% ---------------------------------------------------------------- title
% TODO(author): choose a title. Options (none states a result):
%  (a) [current] Energy and Cost per Query of CPU and GPU Dataframe Engines in the Cloud:
%      A Pre-registered Measurement Protocol and Break-Even Scales
%  (b) When Does a GPU Pay Off for Analytics? Pre-registered Energy and Cost per Query of
%      Polars and Spark on CPU and GPU Cloud Hosts
%  (c) Measure the Meter First: Validated, Scope-Tagged Energy and Cost per Query for CPU and
%      GPU Dataframe Engines on Rented Hosts
\title[Energy and Cost per Query of CPU and GPU Dataframe Engines]{Energy and Cost per Query of
  CPU and GPU Dataframe Engines in the Cloud: A Pre-registered Measurement Protocol and
  Break-Even Scales}

\author{Daniel Ejiofor}
\affiliation{%
  \institution{Boston University}
  \city{Boston}
  \state{MA}
  \country{USA}}
\email{ejiod@bu.edu}

% ---------------------------------------------------------------- abstract
% Result sentences: values and class words from stats.json (Section 6.4 rule; energy in the N1 wording).
\begin{abstract}
Organizations that choose between CPU and GPU engines for analytical data processing rarely
have energy and cost figures per query that were measured on hosts they can rent, with the
meter's scope and resolution stated. We report a pre-registered study of energy and cost per
query for Polars and Apache Spark, each on the CPU and on the GPU (NVIDIA cuDF), on a
TPC-H-derived workload (twelve queries, scale factors 1 to 300), run on NVIDIA A100 and H100
cloud GPU hosts and on a bare-metal CPU host whose RAPL counters are readable. Before the
design was locked we validated each meter on each host class: the counter's scope, the
semantics of the driver's power field, the shortest run it resolves within 2\% ($\tmin$: 1\,s
on the H100, 2\,s on the A100 and on the RAPL path), idle floors with and without an open GPU
context, and clock alignment. The protocol fixes a 10\,s measurement window, reports total and
above-idle energy, tags every value with its scope, and makes cross-class energy claims only
under a stated assumption whose weight is shown by tipping points. Cost is reported as
separate time and price ratios under frozen list prices and three CPU baselines.
At SF300 against the right-sized CPU baseline, the cost ratio of the best CPU engine to the best GPU
engine was
\rtext{x3.h100.baseline.sf300.cost_b1_p_new} on the H100 host (\cls{x3.h100.baseline.sf300.cost_b1_p_new})
and \rtext{x3.a100.baseline.sf300.cost_b1_p_new} on the A100 host (\cls{x3.a100.baseline.sf300.cost_b1_p_new}),
with cost break-even brackets of \stat{break_even.x3.h100.baseline.cost_b1_p_new.text} and
\stat{break_even.x3.a100.baseline.cost_b1_p_new.text}. The energy contrast of the best CPU and GPU engines was
estimable only at the largest scales (SF300, and SF100 for the complex queries), because the fastest GPU
runs fell below the meters' resolution; there it was GPU better at the measured scopes, which supports no
node-level claim: at SF300 the ratio was \rtext{x3.h100.baseline.sf300.energy_total} on the H100 host and
\rtext{x3.a100.baseline.sf300.energy_total} on the A100 host.
Tools and the time-stamped pre-registration are public; the raw records are archived with checksums and
will be deposited.
\end{abstract}

% CCS 2012 concepts verified in ACM's CCS tool (https://dl.acm.org/ccs, read by venue-scout in a browser on
% 2026-11-23): full paths and concept ids from the tree; concept_desc in the tool's "Generate Code" form
% (root~leaf; significance 500 high, 300 medium, 100 low).
%  Information systems > Data management systems > Database administration > Database performance evaluation
%  Hardware > Power and energy > Power estimation and optimization > Enterprise level and data centers power issues
%  General and reference > Cross-computing tools and techniques > Measurement
%  Information systems > Data management systems > Database management system engines > Online analytical processing engines
%  Social and professional topics > Professional topics > Computing industry > Sustainability
\begin{CCSXML}
<ccs2012>
<concept>
<concept_id>10002951.10002952.10003212.10003214</concept_id>
<concept_desc>Information systems~Database performance evaluation</concept_desc>
<concept_significance>500</concept_significance>
</concept>
<concept>
<concept_id>10010583.10010662.10010674.10011724</concept_id>
<concept_desc>Hardware~Enterprise level and data centers power issues</concept_desc>
<concept_significance>300</concept_significance>
</concept>
<concept>
<concept_id>10002944.10011123.10010916</concept_id>
<concept_desc>General and reference~Measurement</concept_desc>
<concept_significance>300</concept_significance>
</concept>
<concept>
<concept_id>10002951.10002952.10003190.10010841</concept_id>
<concept_desc>Information systems~Online analytical processing engines</concept_desc>
<concept_significance>100</concept_significance>
</concept>
<concept>
<concept_id>10003456.10003457.10003458.10010921</concept_id>
<concept_desc>Social and professional topics~Sustainability</concept_desc>
<concept_significance>100</concept_significance>
</concept>
</ccs2012>
\end{CCSXML}
\ccsdesc[500]{Information systems~Database performance evaluation}
\ccsdesc[300]{Hardware~Enterprise level and data centers power issues}
\ccsdesc[300]{General and reference~Measurement}
\ccsdesc[100]{Information systems~Online analytical processing engines}
\ccsdesc[100]{Social and professional topics~Sustainability}

\keywords{energy per query, cost per query, GPU, dataframes, Polars, Apache Spark, NVML, RAPL,
cloud, pre-registration}

\maketitle

% =====================================================================================
\section{Introduction}
\label{sec:intro}

Analytical dataframe engines now run on GPUs as well as CPUs: the Polars GPU engine executes
on NVIDIA cuDF~\cite{cudf_polars_docs}, and Apache Spark~\cite{armbrust2015sparksql} runs SQL
plans on the GPU through the cuDF plugin~\cite{nvidia_cudf_spark_docs}. An organization that
picks an engine and a host type for a workload weighs speed against what each query costs and
how much energy it takes. Speed is what benchmarks usually report. Energy and cost per query
are measured far less often, and when they are, three facts that decide what a number means
are seldom stated: which components the meter covers (its \emph{scope}), the shortest run it
resolves, and how much of the reading is idle draw.

In public clouds this is a structural problem, not an oversight. The GPU instances we could
rent are virtual machines whose guests expose a GPU energy counter but no CPU energy source;
of the hosts we probed, only a bare-metal one exposed readable CPU energy counters (RAPL). A per-query energy figure is
therefore always a figure for a slice of a node, and a CPU-versus-GPU comparison across host
classes compares different slices. For organizations in the United States, where data-center
electricity use is the subject of national reporting~\cite{smith2026datacenter}, the
practical decision is made per workload at hourly list prices that move: one provider in this
study raised three of the four prices relevant to its design by 13--14\% on 2026-11-16, while
the study was being designed. Energy figures that feed procurement or sustainability
reporting need a stated scope, and cost figures need to survive a price change.

We address both with a pre-registered study. Four engines (\polarscpu, \polarsgpu, \sparkcpu,
\sparkgpu) run the TPC-H-derived workload of a prior pre-registered timing
study~\cite{ejiofor2026p1} on two rented GPU hosts (NVIDIA A100 and H100) and on a bare-metal
CPU host with readable RAPL counters; a right-sized CPU virtual machine checks the CPU cost
baseline. The design, including the analysis plan, the classification rule and the scope rules
for what may be claimed, was locked on 2026-11-16 with a public time-stamped
deposit~\citelock{}, after a battery of synthetic-load probes of every meter and before any
timed study session. The analysis code was written after the lock and frozen before any campaign
data was seen. Every later change is a dated amendment reported in Section~\ref{sec:deviations}.

\paragraph{Contributions.}
\begin{enumerate}
\item \textbf{A validated measurement protocol} for per-query energy on rented CPU and GPU
  hosts (Section~\ref{sec:energy}): per host class, the meter's scope, a resolution limit $\tmin$,
  the power field's semantics, idle floors with and without a CUDA context, and clock alignment
  are measured before the lock; a window rule makes sub-second queries measurable; energy is
  reported total and above idle; every value carries a scope tag the analysis code enforces.
\item \textbf{A pre-registered design and analysis plan} (Sections~\ref{sec:design},
  \ref{sec:analysis}): Williams Latin squares, timing-independent integrity gates, a $\pm$10\%
  equivalence margin, a fixed cell classification, and explicit rules for cross-class energy
  claims (assumption \NOne{} and tipping points $\Dstar$), with every deviation disclosed.
\item \textbf{Break-even scales} (Section~\ref{sec:results-breakeven}): the smallest scale
  from which a GPU engine is cheaper per query than the best CPU engine, per suite and host,
  under three CPU baselines and three price tables, and for energy a lower bound. Time and
  price ratios are reported separately, so a reader can re-price the result.
\item \textbf{Open tools and data}: probe tools, harness, gates, price snapshots and raw
  records (Section~\ref{sec:artifacts}).
\end{enumerate}
% One sentence per research question, worded by the rules of Sections 4.6 and 6.4.
Section~\ref{sec:results} reports the answers. RQ1: cross-class energy was estimable only at the larger
scales, and there every cell was GPU better at the measured scopes, which under \NOne{} gives no
node-level direction. RQ2--RQ3: against the right-sized CPU baseline the cost break-even brackets on the
baseline suite are \stat{break_even.x3.h100.baseline.cost_b1_p_new.text} (H100 host) and
\stat{break_even.x3.a100.baseline.cost_b1_p_new.text} (A100 host), every complex-suite scale is GPU better,
and the brackets move with the baseline. RQ4 (descriptive): at SF300 the A100 host took longer than the
H100 host and cost less per query. RQ5: at SF300 a CPU engine cost more on the H100 host than on B1 and
less on the A100 host.

% =====================================================================================
\section{Background and Related Work}
\label{sec:related}

\paragraph{Measuring GPU and CPU energy.}
NVIDIA GPUs expose on-board power sensors through NVML. The library reports a cumulative
energy counter in millijoules since the driver was loaded, and a power reading whose
averaging depends on the architecture~\cite{nvml_api}. Burtscher et al.\ characterized the
built-in sensor of an earlier GPU generation~\cite{burtscher2014k20}, and Yang, Adamek and
Armour report that \texttt{nvidia-smi} samples only about a quarter of runtime on A100 and
H100~\cite{yang2024nvidiasmi}; we therefore difference the cumulative counter and use sampled
power only as a diagnostic. On x86 CPUs, RAPL exposes package and DRAM energy through the
Linux powercap interface~\cite{linux_powercap}. Its update interval limits short
measurements~\cite{hahnel2012rapl}, its behaviour has been characterized at
length~\cite{khan2018rapl}, and its DRAM domain has been validated
separately~\cite{desrochers2016dram}; since the PLATYPUS power side channel~\cite{lipp2021platypus}
the counters are readable by root only. GPU energy has been studied mostly for machine learning~\cite{strubell2019energy,you2023zeus,patel2024llmpower}.
What we add is a per-host-class validation of scope and resolution before measurement, applied
to analytical queries whose runs last from milliseconds to minutes.

\paragraph{Energy in data management.}
The energy efficiency of database servers and clusters has been analyzed on
CPUs~\cite{tsirogiannis2010energy,lang2012cluster,xu2010power}, and TPC result disclosures
have been mined for power consumption~\cite{poess2008energy}; TPC-Energy specifies energy
metrics for TPC benchmarks~\cite{tpcenergy150}.
Nahrstedt et al.\ measured the energy usage and performance of pandas and
Polars~\cite{nahrstedt2024polars}, and Cheng, He and Lau studied energy-efficient query processing
on embedded CPU--GPU hardware~\cite{cheng2015embedded}. In a search of OpenAlex, Crossref and arXiv
we found no measured per-query energy of GPU dataframe engines against CPU engines on rented
server-class hosts with a stated meter scope.
% Search note (2026-11-23; raw responses kept by the writer, not in the repository). Indexes: OpenAlex
% (api.openalex.org/works?search=, title and abstract), Crossref (api.crossref.org/works?query.bibliographic=)
% and arXiv (export.arxiv.org, boolean term queries, cs.DB where given). DBLP was tried and is behind a bot
% check; the ACM DL and IEEE Xplore were not searched directly (their DOIs are in Crossref and OpenAlex).
% Queries (each in OpenAlex and Crossref; 25 and 15 results each): "GPU dataframe energy consumption", "GPU
% database energy efficiency query", "energy per query GPU CPU analytics", "cuDF energy", "RAPIDS
% Accelerator Spark energy", "Spark GPU energy consumption", "Polars energy consumption", "GPU query
% processing power consumption", "GPU SQL analytics energy measurement", "TPC-H energy GPU", "energy
% efficient query processing GPU", "energy consumption GPU join", "energy consumption database
% co-processor", "energy efficiency analytical database GPU CPU comparison", "power measurement
% GPU-accelerated database", "energy GPU accelerated Spark SQL", "energy efficiency dataframe library";
% arXiv: energy AND GPU AND dataframe; energy AND GPU AND query (cs.DB); energy AND Spark AND GPU; energy AND
% RAPIDS; energy AND Polars; energy AND TPC-H AND GPU; power AND GPU AND database (cs.DB). About 800 hits,
% screened by title for energy/power + GPU/accelerator + database/query/dataframe/analytics terms. Relevant:
% Cheng, He, Lau, DaMoN 2015 (embedded CPU-GPU query processing; cited) and Nahrstedt et al., EASE 2024
% (pandas and Polars on the CPU; cited). The other screened titles were GPU performance models, bitmap
% indexes, spatial engines or physics, with no per-query energy of dataframe engines. The claim is worded
% to what this search shows, not as "no one has". Carbon-aware data systems (for the carbon sentence below):
% OpenAlex "carbon-aware query processing", "carbon-aware database system", "carbon aware data analytics
% scheduling", "carbon footprint database query"; relevant: Lechowicz et al., SIGCOMM 2025 (cited there).

\paragraph{Engines, benchmarks and GPU databases.}
Dataframe and analytics engines are compared mostly on speed: the Polars PDS-H
results~\cite{vink2025pdsh}, the db-benchmark~\cite{dbbenchmark}, and studies of scalable
dataframe systems~\cite{petersohn2020dataframes} and embedded engines~\cite{raasveldt2019duckdb}.
GPU database research has surveyed GPU-accelerated systems~\cite{bress2014gpudb,rosenfeld2022heterogeneous},
compared the performance characteristics of GPUs and CPUs for analytics~\cite{shanbhag2020gpucpu},
and studied fast CPU--GPU interconnects~\cite{lutz2020pump}. Lee et al.\ showed how strongly
CPU--GPU speedups depend on the CPU baseline~\cite{lee2010debunking}. The two studies this one
builds on are a pre-registered timing comparison of the same four engines on four GPU
generations~\cite{ejiofor2026p1} and a re-examination of its CPU baseline (library-version
lock-in and core count)~\cite{ejiofor2026p2}, which motivates our three CPU cost baselines. We
follow established benchmarking and reporting
advice~\cite{hoefler2015scientific,raasveldt2018fair,fleming1986geomean}, a Williams design
balanced for first-order carryover~\cite{williams1949latin}, and
pre-registration~\cite{nosek2018prereg}.

\paragraph{Cost per query and carbon.}
Cloud cost per query has been modelled for query processing~\cite{leis2021costoptimal}, cloud
DBMS architectures~\cite{tan2019cloud} and analytics benchmarks~\cite{vanrenen2023cab}.
Carbon-aware computing shifts load in time~\cite{wiesner2021wait,radovanovic2023carbon,lechowicz2025carbon} and
informs data-center design~\cite{acun2023carbonexplorer}, and the footprint of computing
includes embodied as well as operational carbon~\cite{gupta2021chasing}. We report
energy and cost, not carbon: converting our scope-tagged energy to emissions needs grid
intensity and embodied-carbon data that are outside this study.

% =====================================================================================
\section{Study Design}
\label{sec:design}

\subsection{Research questions}
\label{sec:rqs}
The confirmatory questions are fixed in the locked design~\citelock{}. \textbf{RQ1}: energy
per query of a GPU engine against its CPU counterpart, as scope-tagged figures and, under
assumption \NOne{} (Section~\ref{sec:n1}), a node-level direction. \textbf{RQ2}: cost per
query (list price $\times$ time) for the same contrasts against three CPU baselines.
\textbf{RQ3}: the break-even scale, the smallest scale from which a GPU engine is cheaper
than the best CPU engine; for energy only a lower bound under \NOne. \textbf{RQ4}: GPU
generation, A100 against H100 (descriptive; confounded with the host).
\textbf{RQ5}: CPU platform and the idle-GPU tax, i.e.\ CPU engines on the CPU-only host
against the GPU hosts, and the GPU energy burned idle while a CPU engine runs on a GPU host.
``Confirmatory'' means that contrast, metric, baseline, scale grid and decision rule were
fixed in advance (there are no hypothesis tests).

\subsection{Hosts and engines}
\label{sec:hosts}

% Host facts: CPU model and driver of every launch checked against stats.json hosts.<host>.instances
% (2026-11-24); drivers, the C-VM CPU model, vCPU, RAM and storage are record.tex values (the lead read them
% from each launch's campaign.json; both launches of each host are identical).
\begin{table*}[t]
\caption{Hosts. Prices are on-demand list prices in USD per hour from the price table frozen at
the lock (P-NEW; P-OLD in parentheses). ``KVM'' means a verified KVM guest; RAM as seen by the guest.
Disk: \stat{campaign.gpu_hosts.storage} on the GPU hosts, \stat{campaign.gcp_hosts.storage} on the GCP hosts.
GH200 and B200 were part of the design and were not run (amendment A5).}
\label{tab:hosts}
\footnotesize
\setlength{\tabcolsep}{4pt}
\begin{tabular}{@{}l R{0.17\textwidth} R{0.23\textwidth} R{0.09\textwidth} R{0.085\textwidth} R{0.22\textwidth} r@{}}
\toprule
Id & Instance (provider) & CPU as seen by the guest & GPU & vCPU / RAM & Energy meters (scope tag) & USD/h \\
\midrule
\GH{} & \texttt{gpu\_1x\_h100\_sxm5} (Lambda) & Xeon Platinum 8480+, 13 cores / 26 threads, KVM & H100 SXM5 80\,GB & \tstat{campaign.h100.vcpu} / \tstat{campaign.h100.ram_gib}\,GiB & NVML counter, GPU board (\scope{S1-counter}); no CPU energy & 4.89 (4.29) \\
\GA{} & \texttt{gpu\_1x\_a100\_sxm4} (Lambda) & EPYC 7J13, 30 vCPUs, SMT off, KVM & A100 SXM4 40\,GB & \tstat{campaign.a100.vcpu} / \tstat{campaign.a100.ram_gib}\,GiB & NVML counter, GPU board (\scope{S1-counter}); no CPU energy & 1.99 (1.99) \\
\CG{} & \texttt{c3-highmem-192-metal} (GCP) & 2$\times$ Xeon Platinum 8481C, 96 cores / 192 threads, bare metal & none & \tstat{campaign.c-gcp.vcpu} / \tstat{campaign.c-gcp.ram_gib}\,GiB & RAPL package + DRAM per socket (\scope{S2-pinned}, \scope{S2-all}) & 12.701568 \\
\CV{} & \texttt{c3-highmem-44} VM (GCP) & \tstat{campaign.c-vm.cpu_model} & none & \tstat{campaign.c-vm.vcpu} / \tstat{campaign.c-vm.ram_gib}\,GiB & none (time only) & 2.910776 \\
\bottomrule
\end{tabular}
\end{table*}

Table~\ref{tab:hosts} lists the hosts. The GPU hosts are the H100 and A100 instance types of
the prior study~\cite{ejiofor2026p1}, so energy and cost can be read beside its timings; they
are never pooled. On them, all four engines run. The CPU-only host \CG{} is a two-socket
bare-metal machine on which the two CPU engines are pinned to 26 vCPUs of socket~0 (13
physical cores and their SMT siblings, the H100 guest's shape) with memory bound to the
pinned socket's NUMA node (Section~\ref{sec:cpuenergy}). \CV{} checks the right-sized CPU cost baseline (Section~\ref{sec:cost}). The engines and their settings are the prior study's. Every session recorded Polars
\stat{campaign.version.polars}~\cite{polars1442} (CPU and GPU engine) and Spark \stat{campaign.version.spark}
(CPU and GPU), on Python \stat{campaign.version.python}; the GPU engine of Polars is cudf-polars-cu12
\stat{campaign.version.cudf_polars} as pinned in \texttt{requirements-gpu.txt} (installed from the pin; the
version is not recorded per session), and Spark's is the RAPIDS Accelerator jar
\stat{campaign.version.rapids_jar}, the same on all four GPU launches. All run with 26 threads, the strict
\polarsgpu{} session guard (no silent CPU fallback) and the \sparkgpu{} final-plan check. The drivers,
\stat{campaign.h100.driver} on \GH{} and \stat{campaign.a100.driver} on \GA, and host identity
(hypervisor, kernel, storage) are recorded per launch.
Nothing is tuned: no clock locking, power capping or governor change.

\subsection{Workload}
\label{sec:workload}
The workload is the prior study's~\cite{ejiofor2026p1}: a baseline suite of six adapted TPC-H
queries (\q{q01}, \q{q03}, \q{q06}, \q{q10}, \q{q12}, \q{q14}) at SF1, 10, 30, 100 and 300, and
a complex suite of six queries (\q{cx\_q09}, \q{cx\_q13}, \q{cx\_semi\_anti}, \q{cx\_window},
\q{cx\_distinct}, \q{cx\_self\_join}) at SF30, 100 and 300. Inputs are regenerated on each host
with DuckDB's \texttt{tpch} extension, written as Parquet, and must pass a content-fingerprint
gate against the prior study's manifests before round~1; decimal columns are cast to Float64
by every engine; queries are never rewritten. The Polars GPU gaps on \q{cx\_q13} and
\q{cx\_window} are pre-declared coverage outcomes. The scale grid is fixed and never refined
after seeing data. \tpcdisclaimer{} The energy figures here are not TPC-Energy
metrics~\cite{tpcenergy150}.

\subsection{Rounds, order and controls}
\label{sec:rounds}
Each GPU host runs $R = 8$ rounds: the prior study's Williams square for four
engines~\cite{williams1949latin} (rows \texttt{ABDC}, \texttt{BCAD}, \texttt{CDBA},
\texttt{DACB}) for rounds 1--4 and the same rows in reverse order for rounds 5--8, with the
two squares on separate instance launches so that between-instance variation enters the
across-round interval. The engine letters come from
\texttt{random.Random(seed).shuffle} with seed 20261116, the UTC date of the lock. Suite order
alternates by round and scales ascend. A GPU-host round has 32 sessions (4 engines $\times$ 8
suite--scale cells); a \CG{} round has 16, with orders \texttt{AB BA BA AB} and \texttt{BA AB AB
BA} (\sparkcpu{} is A and \polarscpu{} is B, derived from the same seed; only the first square ran, A9). \CV{} runs the
baseline suite at SF10, 30 and 300 for two rounds (\texttt{AB}, \texttt{BA}), time only.
Confirmatory estimates use complete squares only (all of a square's rounds accepted within one
instance): two give $R = 8$, one gives $R = 4$; rounds of an incomplete square are a sensitivity. No round is added or dropped for what its values
look like. One query runs at a time on a whole instance, with no other workload and no foreign
GPU process; at least 150\,s without query load separate consecutive sessions.

% ACM Policy on Authorship, "Use of Artificial Intelligence" (updated 2026-05-14;
% https://www.acm.org/publications/policies/new-acm-policy-on-authorship, read by venue-scout on 2026-11-24):
% AI used to conduct the research "must be described in detail in the methods section of the Work". Facts:
% the lead's list (2026-11-24), design sec. 8 (the design was drafted by Claude via Claude Code; earlier
% project designs with OpenAI Codex) and sec. 11 item 5 (the probes were run by the lead's agents). Names tools,
% not people. Agrees with the Disclosures (submission) and the acknowledgments (camera-ready).
\paragraph{Use of AI tools.} AI coding assistants (Anthropic's Claude Code; OpenAI Codex in the
project's earlier, exploratory work) wrote and tested the measurement harness, the probe tools, the
gates and the campaign automation (launch, drivers, copies, teardown); drafted the pre-registered
design and its amendments; wrote the analysis code, after the lock and frozen before any campaign
data was seen (A7); ran and monitored the probes and the cloud campaign under
\camera{the author's}{the author(s)'} direction; and assisted with writing this paper. \camera{The author}{The author(s)} directed the
study, approved every design and analysis decision and every amendment, reviewed the code and the
text, and \camera{takes}{take} full responsibility. Measurement and analysis are deterministic code
in the study repository; no AI model produces or judges a measured value.

% =====================================================================================
\section{Energy Measurement Protocol}
\label{sec:energy}

\subsection{Meters, scopes and scope tags}
\label{sec:scopes}

\paragraph{GPU energy (\scope{S1}).}
The primary GPU source is NVML's \texttt{nvmlDeviceGetTotalEnergyConsumption}, a cumulative
millijoule counter~\cite{nvml_api}, read through \texttt{ctypes} because \texttt{nvidia-smi
--query-gpu} cannot return it, and differenced per run. The worker stores UTC and
\texttt{CLOCK\_MONOTONIC} stamps just outside the prior study's timer and reads the counter just outside
them; read latency and overhang are recorded per run. Before the lock we checked on each
host class that the counter advances (update period about 100\,ms on the H100) and has the
scope we claim for it, by comparing it with the 20\,ms integral of an instantaneous-power field
(Table~\ref{tab:validation} in Appendix~\ref{app:validation}). On the H100 and A100 the counter
agrees with the GPU-only
\texttt{power.draw.instant} integral within 0.3\%, the drivers expose no module-level fields, and
the counter is the GPU board (GPU, HBM, voltage regulators; the host is excluded): tag
\scope{S1-counter}. The check matters: on the GH200 the counter matched the module field (0.998),
not the GPU-only one (1.14--1.55); it covers the whole module (Appendix~\ref{app:notrun}).

\paragraph{CPU energy (\scope{S2}) and where it is missing.}
\label{sec:cpuenergy}
On \CG{} the source is RAPL through powercap~\cite{linux_powercap}: one \texttt{package} and one
\texttt{dram} zone per socket, all four advancing. The primary scope \scope{S2-pinned} is package
plus DRAM of socket~0, which hosts the engine threads; \scope{S2-all} adds the idle socket~1,
which is reported, not attributed. A 200\,ms sampler reads the counters (unwrapping with
\texttt{max\_energy\_range\_uj}); a window that contains a counter decrease or a gap above 30\,s
is invalid, and per-run energy is the series interpolated linearly at the run's monotonic stamps.
Because \texttt{energy\_uj} is root-only~\cite{lipp2021platypus}, read access was granted on this
single-tenant host by a setup flag, recorded. The RAPL path and the pinning were verified before the lock (Appendix~\ref{app:validation}). DRAM attribution to the pinned socket is therefore sound.

The Lambda GPU hosts are KVM guests with no CPU energy source: no powercap zones, no MSR device,
no BMC, and a \texttt{perf} \texttt{energy-psys} event that counts 0\,J (H100; a source check on
the 1$\times$ A100 found none either). CPU and DRAM energy of a GPU host is therefore not measured.
No model value enters a confirmatory table: every energy value carries one of the tags
\scope{S1-counter}, \scope{S1-sampled}, \scope{S1-module}, \scope{S2-pinned}, \scope{S2-all},
\scope{S3} (whole node by BMC; no host had one) or \scope{none}, and the analysis code raises an
error on an untagged value. The design's exploratory host-CPU energy model (MODEL, design
\S\S2.4--2.5) is not computed (A11); no model value appears in this paper.

\paragraph{The \texttt{power.draw} field and the sampled cross-check.}
The default \texttt{nvidia-smi} power field means different things on different GPUs
(Table~\ref{tab:validation}): on the H100 it equals \texttt{power.draw.average} in every sample
(a trailing 1\,s average; correlation with the instantaneous field 0.03), on the A100 it is
instantaneous, consistent with the NVML documentation~\cite{nvml_api}. Energy integrated from it
is thus not like-for-like across GPUs. We record \texttt{.average} and \texttt{.instant} in a
200\,ms stream and use its trapezoid integral only as a diagnostic cross-check on windows of at
least 10\,s.

\subsection{Resolution: the window rule and \texorpdfstring{$\tmin$}{t\_min}}
\label{sec:resolution}
A counter that updates every 100\,ms (NVML) or a series sampled every 200\,ms (RAPL) leaves each
end of a run uncertain by up to one period, which is large for a 50\,ms query. A measured
\emph{trial} is therefore a window of back-to-back runs of one query: runs continue until their
summed timed time reaches 10\,s or 50 runs have been made ($n$ = runs made). Window energy is the
sum of per-run counter differences, gaps, cleanup and plan explanation excluded; per-run energy
and time are the window values divided by $n$. Each query gets three trials and the median is the
round's value. $n$ depends only on elapsed time, never on energy.

How short a run can be and still be measured was itself measured before the lock (probe P6). On
each host class, synthetic pulses of known length (0.1 to 5\,s, plus isolated pulses up to
10\,s) were measured by the campaign's own energy path and compared with the 20\,ms integral of
the instantaneous field of the counter's scope (on \CG: the 200\,ms interpolated series against
direct \texttt{energy\_uj} reads at the same stamps), for single pulses and for windows built by
the window rule. The pre-stated rule is: $\tmin$ is the smallest per-run length whose median
absolute error is at most 2\% and every longer length's also is, in windows mode, with the
locked 1\,s default as a floor, $\tmin(\text{host}) = \max(1\,\text{s},\ \text{measured})$. This
gives 1\,s on the H100 (measured 0.5\,s), 2\,s on the A100 and 2\,s on \CG{}
(Table~\ref{tab:validation}). A cell whose median trial has a per-run time below its host's
$\tmin$, or a window below 10\,s, is reported as \emph{below measurement resolution}, never as a
number (gate G9); the cell's time and cost stay. The limit was not tuned to observed values: the A100's 1.9\% at 2\,s is close
to the limit and is recorded as measured.

The errors are not always random. On the H100 their signs are mixed, so for $n$ runs they grow
like $\sqrt{n}$; on the A100 every window was positive at every length (about 4--12\,J per run),
and on \CG{} every window and pulse was negative (7.0--8.9\,J per run at a spin load of about
340\,W for package~0 and its DRAM, consistent with the 200\,ms sampling smearing each power step;
not tested). One-sided errors do not average out, which is why $\tmin$ is taken from windows
mode, the mode the study uses. The single-run $\tmin$ never binds (Appendix~\ref{app:validation}).

\subsection{Idle, both ways}
\label{sec:idle}
Each session is bracketed by two 30\,s idle windows with the engine ready and no query running,
each preceded by at least 60\,s without load. $\Pidle$ per scope is the mean of the two windows'
mean power. Every energy value is reported \emph{total}, $E$, the counter difference over the
timed runs (primary: what the machine draws), and \emph{above idle}, $\Eabove = E - \Pidle\, t$,
negative values reported as they are. For a CPU engine on a GPU host the GPU-board above-idle
energy is zero by construction and not reported.

Two probe results shape how idle numbers can be read (Table~\ref{tab:validation}). First, an
open CUDA context with no kernel running raises idle board power by 68\,W on the H100 (SM clock
1980 against 345\,MHz) but by 5.6\,W on the A100, so a GPU-engine session's $\Pidle$ and the idle
GPU draw under a CPU-engine session are different quantities; which state a session's windows
measured is read from its recorded clocks. Second, a 26-process CPU burn moved the A100 counter
by nothing and raised the H100 board counter by 3--5\% together with GPU temperature, a GPU-side
effect: host CPU load can move the ``idle GPU'' energy of a CPU-engine run. The measurement processes run in every session, so their draw is inside $\Pidle$ and $E$
(idle stability and overhead: Appendix~\ref{app:validation}).

\subsection{Clock alignment}
\label{sec:clock}
Counter energy needs no cross-process clock alignment (the worker stamps and reads in one
process), but the 200\,ms \texttt{nvidia-smi} stream does. Gate G8 checks it from onset times
only, never magnitudes: the first stream sample with GPU utilization of at least 10\% must
follow the worker's window start with a median lag within $\pm 0.4$\,s and a first-third
versus last-third drift of at most 0.2\,s, using GPU-engine windows of at least 5\,s, or else 3\,s
synthetic pulses run before the worker starts and after it exits. A failed G8 disables only that
session's stream values; counter energy and the round are untouched. Before the lock the G8
probe passed on every host with no onset missed, the A100 marginally (median 0.389\,s; 8 of 20
lags above 0.4\,s; Table~\ref{tab:validation}). Before any study data, the smoke tests led to a stricter onset reading (A3) and to the stated
expectation that many sessions may carry a G8 flag, because a GPU engine's onset lag includes its
start-up latency (A4); flags are reported per host and engine.

\subsection{What may be compared: assumption \NOne{} and tipping points}
\label{sec:n1}
Within a host class, energy contrasts use one scope. Across classes there is no common scope: a
GPU-engine run is measured on the GPU board ($E_G$, \scope{S1-counter}) and a CPU-engine run on
\CG{} as package plus DRAM ($E_C$, \scope{S2-pinned}), each a lower bound on its node's energy.
Write node energy as measured plus unmeasured: $E_G + U_G$ (host CPU, DRAM and platform of the
GPU node) and $E_C + U_C$ (the platform of the CPU node). The design adopts, as a labelled
assumption that is not measured or tested,
\begin{equation}
\NOne:\quad U_G \ge U_C,
\end{equation}
plausible because the GPU node adds a host CPU and DRAM running the same engine process to the
same kind of platform, but it may fail. Under \NOne{} the direction ``the GPU engine uses
\emph{more} node energy'' survives when a cell is CPU better at the measured scopes, and the
measured-scope break-even is a \emph{lower bound} on the node-level energy break-even. ``The GPU
engine uses less or equal energy'', an energy break-even scale itself, the magnitude of any
cross-class ratio and ``no difference'' across classes do not survive. Time and cost claims need
no assumption. For every cross-class cell we report the tipping point
\begin{equation}
\Dstar = E_C - E_G,
\end{equation}
the unmeasured excess $U_G - U_C$ at which node energies would be equal, from per-round suite
totals (mean and 95\% interval, in Wh per suite run and in W over the GPU suite time), beside
the measured idle power of one \CG{} socket and of the idle GPU for scale. If $\Dstar < 0$, \NOne{}
already implies that the GPU node uses more energy; if $\Dstar > 0$ we make no claim and $\Dstar$
says how large the GPU node's unmeasured excess would have to be to erase its measured
advantage. Hence the standing limitation of this study: on the x86 GPU hosts, ``a GPU engine
uses less energy than its CPU counterpart'' cannot be claimed at any scale.

% =====================================================================================
\section{Cost Model}
\label{sec:cost}

The cost of a run is $c = p\,t/3600$, with $p$ the on-demand list price in USD per hour of the
instance rented and $t$ the timed seconds, accrued per second; energy, carbon, storage, egress
and billing minimums are not in per-query cost. Session cost (worker launch to process exit,
less the idle windows) and a cold-job cost $(p/3600)(t_{\mathrm{ready}} + t)$ are secondary in the
design and are not computed (A11).

\paragraph{Price tables.}
\textbf{P-NEW} (primary) is frozen at the lock: Lambda prices effective 2026-11-16 00:00 UTC
and Google's official list prices fetched 2026-11-11 (Table~\ref{tab:hosts}). \textbf{P-OLD}
holds Lambda's prices before that date (H100 4.29). \textbf{P-BILLED} is the price the provider
charged, from each launch record, with the invoice reconciled against hours $\times$ price and
any gap above 2\% reported. Every session is costed under all three whatever was billed and
whatever credits funded the bill, so cost ratios under P-NEW depend only on measured time and a
table fixed before any outcome existed.

\paragraph{CPU baselines.}
Following the lesson that the CPU baseline drives the conclusion~\cite{lee2010debunking,ejiofor2026p2},
a CPU engine is never charged for an idle GPU in the primary baseline. \textbf{B1} (primary) is
a right-sized CPU VM, \texttt{c3-highmem-44} at 2.910776\,USD/h (the smallest C3 shape with
$\ge$26 vCPUs and $\ge$200\,GiB, the RAM of the smallest GPU host), applied to times measured on
\CG{} pinned to 26 vCPUs. \CV{} checks this: B1 holds if the geometric-mean time ratio
VM/metal is in $[0.90, 1.10]$ and no CPU-engine query fails for resource exhaustion at SF300;
otherwise B1 uses the \CV{} times where they exist and the gap is reported.
\textbf{B2} is \CG{} as rented (an inflated upper bound) and \textbf{B3} the same CPU engine on
the same GPU host (the idle-GPU tax).

\paragraph{Factorised reporting.}
A GPU host is cheaper per query if and only if $p_G/p_C < t_C/t_G$. We report the time ratio and the price ratio separately. The
price ratios, fixed at the lock, are the speedups a GPU engine needs to cost less: against B1,
1.680 on the H100 (1.474 under P-OLD) and 0.684 on the A100; against B2, 0.385 (0.338) and 0.157;
under B3, 1. A reader with other prices can re-price every result with their own ratio.

% =====================================================================================
\section{Gates and Analysis Plan}
\label{sec:analysis}

\subsection{Integrity and telemetry gates}
\label{sec:gates}
Gates are pass/fail and never based on engine comparisons (Table~\ref{tab:gates} in
Appendix~\ref{app:gates}). G1--G5 are
the prior study's: every session exits without timeout; every engine's output agrees with
\polarscpu{} (relative $10^{-8}$, absolute $10^{-6}$, integers exact); the strict GPU guard; the
Spark GPU plan check; and verified garbage-collection barriers in the complex suite. Rounds fail
on G1--G6 and G12 and are repeated once; a launch whose round fails twice, or whose instance is
lost, is relaunched once on a fresh instance, restarting its square, and the lost launch's rounds
enter no estimate (amendment A1). G6 requires each session's counters to match the
probe-declared sources and to advance; G7 checks the sampled stream's coverage, G8 clock alignment
(Section~\ref{sec:clock}), G9 resolution, G10 idle stability (10\%), G11 thermal and power state,
and G12 exclusive use. G7--G11 only set eligibility flags. On GPU hosts a
CUDA-readiness check runs before any data or engine work (A4; Section~\ref{sec:deviations}).

\subsection{Quantities and contrasts}
\label{sec:contrasts}
For each host, round, suite, scale, query and engine, $m$ is the median over three trial windows
of per-run time, energy (total; above-idle alongside) and cost. A contrast is
$\rho = m_{\mathrm{ref}}/m_{\mathrm{engine}}$, so $\rho > 1$ means the engine needs less time,
energy or cost than the reference. Per host, suite, scale and metric with a valid scope or price:
\textbf{X1} \polarscpu{}/\polarsgpu{}; \textbf{X2} \sparkcpu{}/\sparkgpu{}; \textbf{X3} best CPU
engine / best GPU engine per query and round among engines that ran the query (strict variant:
queries both GPU engines ran); if an engine that ran the query has no value for the metric (for
example, below resolution), that query's X3 value is not estimable. The CPU reference is B1 for cost and \scope{S2-pinned} on \CG{} for
energy, so energy X1--X3 are cross-class and read under \NOne{} with $\Dstar$. \textbf{X4}: per GPU
engine, A100 over H100 (descriptive, confounded with CPU, RAM, link and driver). \textbf{X5}: CPU
engines on a GPU host over \CG{}, time and cost only. \textbf{X6}: the idle-GPU tax, B3/B1 as price
ratio $\times$ time ratio, and the \scope{S1} idle-GPU energy of CPU-engine runs in Wh per suite
run. The design's exploratory analyses (EX1--EX3, among them a re-analysis of the prior study's
archived GPU power~\cite{ejiofor2026rawp1}) are not computed (A11).

\subsection{Estimation}
\label{sec:estimation}
Per query, the log-ratio is computed per round; the estimate is the geometric mean over rounds
and the 95\% interval is a $t$-interval on the mean log-ratio with $R-1$ degrees of freedom
($t = 2.365$ for $R = 8$, $3.182$ for $R = 4$), back-transformed; the across-round range is also
reported. Per suite and scale, the per-query log-ratios of queries both engines supported are
averaged within each round and the same interval is taken across rounds (round as block).
Cross-host contrasts (cross-class energy, B1 and B2 cost, X5, X6) pair the two hosts' rounds by
round number, a convention because the rounds are independent runs; an unpaired Welch estimate
is shown beside each such cell as a labelled sensitivity. Suite totals in Wh and USD per suite
run are reported beside geometric means, never mixed with them.
There are no hypothesis tests and no $p$-values; intervals are marginal, with no multiplicity
correction and no power calculation (between-round energy variance was unknown). Ineligible
cells (coverage outcome, below resolution, scope absent, invalid session) go to a coverage
matrix and are not imputed.

\subsection{Classification and break-even scale}
\label{sec:classification}
Each (contrast, host, suite, scale) cell is classified from its 95\% interval of $\rho$:
\emph{GPU better} if it lies entirely above 1, \emph{CPU better} if entirely below 1, \emph{no
difference} if entirely inside $[1/1.10,\ 1.10]$ (a practical margin approved before data), and
\emph{inconclusive} otherwise, checked in this order (an interval wholly above 1 but inside the
margin is GPU better and flagged as within the margin); an interval that merely contains 1 is never
``no difference''.
The \emph{break-even scale} of a contrast, per suite, host and metric (energy total and above
idle; cost under B1, B2, B3, each under P-NEW, P-OLD and P-BILLED) is the smallest grid scale $s$
such that every grid scale $\ge s$ is GPU better, reported as the bracket (previous scale, $s$];
``$\le$SF1'' and ``$>$SF300'' (the top scale has a class other than GPU better) are censored
outcomes, and a scale that is not estimable is named in the outcome (``$\le$SF10 (SF1 not estimable)''; ``not determinable (SF300 not estimable)'').
For cross-class energy it is computed on the
measured scopes and labelled a lower bound on the node-level break-even, under \NOne. No scale is
added to refine a bracket.

\subsection{Blinding}
\label{sec:blinding}
Until every round has completed or been stopped by rule, no engine-versus-engine or
host-versus-host energy, power or cost aggregate is computed or viewed; only QA views (gate
outcomes, sample counts, idle stability, coverage, read latencies) are allowed, through an
allow-list view of the logs (A2). The analysis code (A7) refuses the full analysis until
every planned launch is complete, stopped or declared ended, and before that reads only the gate
files of the QA view. Stopping, repeating or relaunching is never based on observed
speed, energy or cost.

% =====================================================================================
\section{Deviations from the Pre-registration}
\label{sec:deviations}

The design was locked on 2026-11-16 (annotated tag, commit \lockcommit{}, deposit \lockdoi{});
what had been seen before is disclosed in the design (synthetic-load probe results, the prior
study's published timings and GPU power summaries). Sections 1--11 of the design were never
edited afterwards. Eleven dated amendments followed, each stating the outcome data seen when it was
made (``seen''): A1 (lock-day price snapshot, lost-launch policy, SF1 preflight, the \CG{}
allow-list as pending, record completeness, start order, funding, a synthetic-data guard; seen:
none), A2 (B200 admission, the bring-up view; seen: none beyond synthetic probe load), A3 (smoke
test 1, the G8 onset reading; seen: none beyond synthetic load), A4 (smoke tests 2 and 3, the
CUDA-readiness gate, the expected G8 flags; seen: none beyond synthetic load), A5 (GH200 and B200
not run, GCP funding, start order; seen: none), A6 (the \CG{} allow-list; seen: none), A7 (the
analysis code; seen: none), A8 (a gate defect on the CPU-only hosts and their relaunch; seen: none)
A9 (\CG{} launch~2 not run, for time; seen: none), A10 (the analysis command and flags, fixed before
unblinding; seen: none) and A11 (secondary and exploratory analyses not computed; seen: none). None changes a threshold, estimator, margin or primary definition.
Appendix~\ref{app:amendments} is the full pre-registration record (Table~\ref{tab:deviations}).

Six deviations affect what the paper can report. First, contrary to \S5.4 of the locked
design~\citelock{}, no analysis code existed at the lock: the analysis code was written after the
lock and frozen before any campaign data was seen (A7 item 1). It was written by a team member
without access to the campaign data, validated only on synthetic campaigns with injected known
effects, and independently reviewed; amendment A7 (\stat{amend.a7.date}, commit
\camera{\stat{amend.a7.commit}}{\anonmark}) records the freeze commit
(\camera{\stat{amend.a7.freeze_commit}}{\anonmark}) and the code's sha256. On synthetic campaigns all
1,056 X1--X3 estimates lay within sampling error of the injected truth, and the 95\% intervals
covered it in 3,079 of 3,240 cases (0.950) over 12 replicate seeds (A7 item 3). The
pre-registration therefore fixes the analysis rules of Section~\ref{sec:analysis} but not, in
advance, the code that applies them; the readings of the design that the code applies are listed in
Appendix~\ref{app:readings}.
Second, GH200 and B200 were dropped (A5.1) by
the author's decision on 2026-11-23, on grounds of time and scarce capacity, before any campaign
value existed: RQ4 covers two GPU generations instead of four, and the planned same-host,
module-scope CPU-versus-GPU energy comparison on the GH200, the one comparison in the design that
needed no \NOne, is not available. Their pre-lock probes remain as method validation
(Appendix~\ref{app:notrun}). Third, the relaunch rule (A1.2, Section~\ref{sec:gates}) decides which instances form a square;
the lost launch's rounds are reported, never pooled. Fourth, a GPU host that fails the
CUDA-readiness check before round~1 is replaced without using that relaunch (A4.2, the lead's
reading). Fifth, a round-gate defect on the CPU-only hosts (A8: G1 required an idle-GPU record
written only for GPU-engine jobs) stopped both of their launch~1s; after the fix each was relaunched
once under A1.2, which changed the instances that form the CPU-only squares. These three depend only
on technical failures, never on a value; lost and replaced launches are counted in Table~\ref{tab:qa}.
Sixth, \CG{} launch~2 was not run (A9: the author's decision, by the clock, before any result was seen),
so \CG's confirmatory estimates come from one square and every cell paired with \CG{} uses rounds 1--4.
The design's secondary session energy, session cost and cold-job cost and its exploratory EX1--EX3,
including the MODEL host-CPU estimate, are not computed (A11); none is confirmatory, and the deposited raw
data allow them later as labelled exploratory analyses.

% =====================================================================================
\section{Results}
\label{sec:results}
% Results skeleton (2026-11-24): every number is a stats key; every judgment is the analysis's own class
% word, printed by \cls / \ecls / \rtext / \dstar (macros.tex), so the prose holds whichever way a result
% falls, including "not estimable" and censored brackets. Read against the output on 2026-11-24; the
% interpretation is in Section 9 only. A host with no complete square would be reported descriptively.

% Full-width results floats are defined here, at the start of the section, so that they can
% float onto the results pages (a table* lands at the top of a page after the one it is defined on).
\begin{figure*}[t]
\centering
% Drawn by tools/make_figures.py from stats.json (display rules in its header); placeholder sized as the figure.
\studyfig{f01_energy_forest}{\textwidth}{3.4in}
\caption{Energy ratio $\rho$ of X1, X2 and X3 at the measured scopes (\scope{S2-pinned} on \CG{} over
\scope{S1-counter} on the GPU host; log scale; whiskers: 95\% interval) per host, suite and scale; top row
total energy, bottom row above idle. Marker shape gives the class; hollow: one round, no interval;
cross: not estimable; shaded: the $\pm$10\% band.}
\label{fig:energy}
\end{figure*}

\begin{table*}[t]
\caption{Classification of X3 (best CPU engine / best GPU engine) per host, suite and
scale; cost under B1, B2, B3 (P-NEW) and energy at the measured scopes (total). Cell: estimate
of $\rho$ and class (G: GPU better; C: CPU better; =: no difference; ?: inconclusive; --: not
estimable, e.g.\ below resolution); 95\% intervals in Appendix~\ref{app:tables}. Keys
\texttt{x3.<host>.<suite>.sf<scale>.\{cost\_b1\_p\_new, cost\_b2\_p\_new, cost\_b3\_p\_new,
energy\_total\}}; X1, X2, P-OLD and P-BILLED in Appendix~\ref{app:tables}.}
\label{tab:x3}
\footnotesize
\setlength{\tabcolsep}{4pt}
\begin{tabular}{@{}ll*{8}{c}@{}}
\toprule
& & \multicolumn{5}{c}{baseline suite} & \multicolumn{3}{c}{complex suite} \\
\cmidrule(lr){3-7}\cmidrule(l){8-10}
Host & Metric & SF1 & SF10 & SF30 & SF100 & SF300 & SF30 & SF100 & SF300 \\
\midrule
\multirow{4}{*}{\GH} & cost B1 & \xcell{x3.h100.baseline.sf1.cost_b1_p_new} & \xcell{x3.h100.baseline.sf10.cost_b1_p_new} & \xcell{x3.h100.baseline.sf30.cost_b1_p_new} & \xcell{x3.h100.baseline.sf100.cost_b1_p_new} & \xcell{x3.h100.baseline.sf300.cost_b1_p_new} & \xcell{x3.h100.complex.sf30.cost_b1_p_new} & \xcell{x3.h100.complex.sf100.cost_b1_p_new} & \xcell{x3.h100.complex.sf300.cost_b1_p_new} \\
 & cost B2 & \xcell{x3.h100.baseline.sf1.cost_b2_p_new} & \xcell{x3.h100.baseline.sf10.cost_b2_p_new} & \xcell{x3.h100.baseline.sf30.cost_b2_p_new} & \xcell{x3.h100.baseline.sf100.cost_b2_p_new} & \xcell{x3.h100.baseline.sf300.cost_b2_p_new} & \xcell{x3.h100.complex.sf30.cost_b2_p_new} & \xcell{x3.h100.complex.sf100.cost_b2_p_new} & \xcell{x3.h100.complex.sf300.cost_b2_p_new} \\
 & cost B3 & \xcell{x3.h100.baseline.sf1.cost_b3_p_new} & \xcell{x3.h100.baseline.sf10.cost_b3_p_new} & \xcell{x3.h100.baseline.sf30.cost_b3_p_new} & \xcell{x3.h100.baseline.sf100.cost_b3_p_new} & \xcell{x3.h100.baseline.sf300.cost_b3_p_new} & \xcell{x3.h100.complex.sf30.cost_b3_p_new} & \xcell{x3.h100.complex.sf100.cost_b3_p_new} & \xcell{x3.h100.complex.sf300.cost_b3_p_new} \\
 & energy & \xcell{x3.h100.baseline.sf1.energy_total} & \xcell{x3.h100.baseline.sf10.energy_total} & \xcell{x3.h100.baseline.sf30.energy_total} & \xcell{x3.h100.baseline.sf100.energy_total} & \xcell{x3.h100.baseline.sf300.energy_total} & \xcell{x3.h100.complex.sf30.energy_total} & \xcell{x3.h100.complex.sf100.energy_total} & \xcell{x3.h100.complex.sf300.energy_total} \\
\midrule
\multirow{4}{*}{\GA} & cost B1 & \xcell{x3.a100.baseline.sf1.cost_b1_p_new} & \xcell{x3.a100.baseline.sf10.cost_b1_p_new} & \xcell{x3.a100.baseline.sf30.cost_b1_p_new} & \xcell{x3.a100.baseline.sf100.cost_b1_p_new} & \xcell{x3.a100.baseline.sf300.cost_b1_p_new} & \xcell{x3.a100.complex.sf30.cost_b1_p_new} & \xcell{x3.a100.complex.sf100.cost_b1_p_new} & \xcell{x3.a100.complex.sf300.cost_b1_p_new} \\
 & cost B2 & \xcell{x3.a100.baseline.sf1.cost_b2_p_new} & \xcell{x3.a100.baseline.sf10.cost_b2_p_new} & \xcell{x3.a100.baseline.sf30.cost_b2_p_new} & \xcell{x3.a100.baseline.sf100.cost_b2_p_new} & \xcell{x3.a100.baseline.sf300.cost_b2_p_new} & \xcell{x3.a100.complex.sf30.cost_b2_p_new} & \xcell{x3.a100.complex.sf100.cost_b2_p_new} & \xcell{x3.a100.complex.sf300.cost_b2_p_new} \\
 & cost B3 & \xcell{x3.a100.baseline.sf1.cost_b3_p_new} & \xcell{x3.a100.baseline.sf10.cost_b3_p_new} & \xcell{x3.a100.baseline.sf30.cost_b3_p_new} & \xcell{x3.a100.baseline.sf100.cost_b3_p_new} & \xcell{x3.a100.baseline.sf300.cost_b3_p_new} & \xcell{x3.a100.complex.sf30.cost_b3_p_new} & \xcell{x3.a100.complex.sf100.cost_b3_p_new} & \xcell{x3.a100.complex.sf300.cost_b3_p_new} \\
 & energy & \xcell{x3.a100.baseline.sf1.energy_total} & \xcell{x3.a100.baseline.sf10.energy_total} & \xcell{x3.a100.baseline.sf30.energy_total} & \xcell{x3.a100.baseline.sf100.energy_total} & \xcell{x3.a100.baseline.sf300.energy_total} & \xcell{x3.a100.complex.sf30.energy_total} & \xcell{x3.a100.complex.sf100.energy_total} & \xcell{x3.a100.complex.sf300.energy_total} \\
\bottomrule
\end{tabular}
\end{table*}


\subsection{Campaign record and data quality}
\label{sec:results-qa}

\begin{table}[t]
% Keys (analyze_paper3.py, format frozen 2026-11-23): hosts.<host>.status; counts.<host>.{squares_complete,
% squares_planned, rounds_confirmatory, launches_lost_replaced, round_repeats, sessions_analysed,
% g9.cells_below_resolution, coverage_outcomes.total}; c_vm_check.status; qa.g8.<host>.<engine>.{flagged,
% sessions}; qa.crosscheck.<host>.median ((trapezoid - counter) / counter, windows >= 10 s, sessions with
% clean G7 and G8); qa.read_latency_and_overhang.<host>.polars_gpu.read_latency_ms.median.
\caption{Campaign record and QA per host. G8: sessions flagged (failed or not evaluable) of
sessions. Cross-check: median relative difference of the sampled \texttt{nvidia-smi} integral from
the counter (windows of at least 10\,s; sessions with clean G7 and G8).}
\label{tab:qa}
\footnotesize
\setlength{\tabcolsep}{3pt}
\begin{tabular}{@{}lcccc@{}}
\toprule
Item & \GH & \GA & \CG & \CV \\
\midrule
Status & \tstat{hosts.h100.status} & \tstat{hosts.a100.status} & \tstat{hosts.c-gcp.status} & \tstat{hosts.c-vm.status} \\
Complete squares & \tstat{counts.h100.squares_complete}/\tstat{counts.h100.squares_planned} & \tstat{counts.a100.squares_complete}/\tstat{counts.a100.squares_planned} & \tstat{counts.c-gcp.squares_complete}/\tstat{counts.c-gcp.squares_planned} & -- \\
Confirmatory rounds & \tstat{counts.h100.rounds_confirmatory} & \tstat{counts.a100.rounds_confirmatory} & \tstat{counts.c-gcp.rounds_confirmatory} & \tstat{counts.c-vm.rounds_confirmatory} \\
Launches lost or replaced & \tstat{counts.h100.launches_lost_replaced} & \tstat{counts.a100.launches_lost_replaced} & \tstat{counts.c-gcp.launches_lost_replaced} & \tstat{counts.c-vm.launches_lost_replaced} \\
Round repeats & \tstat{counts.h100.round_repeats} & \tstat{counts.a100.round_repeats} & \tstat{counts.c-gcp.round_repeats} & \tstat{counts.c-vm.round_repeats} \\
Sessions analysed & \tstat{counts.h100.sessions_analysed} & \tstat{counts.a100.sessions_analysed} & \tstat{counts.c-gcp.sessions_analysed} & \tstat{counts.c-vm.sessions_analysed} \\
Cells below resolution (G9) & \tstat{counts.h100.g9.cells_below_resolution} & \tstat{counts.a100.g9.cells_below_resolution} & \tstat{counts.c-gcp.g9.cells_below_resolution} & -- \\
Coverage outcomes (cells) & \tstat{counts.h100.coverage_outcomes.total} & \tstat{counts.a100.coverage_outcomes.total} & \tstat{counts.c-gcp.coverage_outcomes.total} & \tstat{counts.c-vm.coverage_outcomes.total} \\
B1 time check & -- & -- & -- & \tstat{c_vm_check.status} \\
\midrule
G8, \polarscpu & \tstat{qa.g8.h100.polars.flagged}/\tstat{qa.g8.h100.polars.sessions} & \tstat{qa.g8.a100.polars.flagged}/\tstat{qa.g8.a100.polars.sessions} & -- & -- \\
G8, \polarsgpu & \tstat{qa.g8.h100.polars_gpu.flagged}/\tstat{qa.g8.h100.polars_gpu.sessions} & \tstat{qa.g8.a100.polars_gpu.flagged}/\tstat{qa.g8.a100.polars_gpu.sessions} & -- & -- \\
G8, \sparkcpu & \tstat{qa.g8.h100.spark_cpu.flagged}/\tstat{qa.g8.h100.spark_cpu.sessions} & \tstat{qa.g8.a100.spark_cpu.flagged}/\tstat{qa.g8.a100.spark_cpu.sessions} & -- & -- \\
G8, \sparkgpu & \tstat{qa.g8.h100.spark_gpu.flagged}/\tstat{qa.g8.h100.spark_gpu.sessions} & \tstat{qa.g8.a100.spark_gpu.flagged}/\tstat{qa.g8.a100.spark_gpu.sessions} & -- & -- \\
Cross-check, median & \tstat{qa.crosscheck.h100.median} & \tstat{qa.crosscheck.a100.median} & -- & -- \\
Read latency ms, \polarsgpu & \tstat{qa.read_latency_and_overhang.h100.polars_gpu.read_latency_ms.median} & \tstat{qa.read_latency_and_overhang.a100.polars_gpu.read_latency_ms.median} & -- & -- \\
\bottomrule
\end{tabular}
\end{table}

Sessions ran from \stat{campaign.first_session_utc} to \stat{campaign.last_session_utc}; Table~\ref{tab:qa}
gives the record. Confirmatory estimates use complete squares only (Section~\ref{sec:rounds}):
\stat{counts.h100.squares_complete} of \stat{counts.h100.squares_planned} on \GH{}
($R={}$\stat{counts.h100.rounds_confirmatory}), \stat{counts.a100.squares_complete} of
\stat{counts.a100.squares_planned} on \GA{} ($R={}$\stat{counts.a100.rounds_confirmatory}) and
\stat{counts.c-gcp.squares_complete} of \stat{counts.c-gcp.squares_planned} on \CG{}
($R={}$\stat{counts.c-gcp.rounds_confirmatory}). A cross-host contrast pairs rounds by round number
(Appendix~\ref{app:readings}), so its degrees of freedom, reported with each interval, follow the
rounds both hosts have. Launches were relaunched or replaced only after technical failures (A1.2,
A4.2, A8); \CG{} launch~2 was not run (A9). G9 marked \stat{counts.h100.g9.cells_below_resolution} of
\stat{counts.h100.cells_analysed} cells below resolution on \GH{} (all with a per-run time below
$\tmin$; \stat{counts.h100.g9.window_below_10s} also with a window below 10\,s),
\stat{counts.a100.g9.cells_below_resolution} of \stat{counts.a100.cells_analysed} on \GA{} and
\stat{counts.c-gcp.g9.cells_below_resolution} of \stat{counts.c-gcp.cells_analysed} on \CG. Their energy is
not estimable (their time and cost stay), and an X3 cell is not estimable when any engine that ran lacks a
value (Appendix~\ref{app:readings}), so cross-class energy exists only at the larger scales: X1 and X3 at
SF300 on the baseline suite and at SF100--SF300 on the complex suite, X2 from SF30 on (Table~\ref{tab:x3},
Table~\ref{tab:x12}). Coverage outcomes affected \stat{counts.h100.coverage_outcomes.total} cells on
\GH{} and \stat{counts.a100.coverage_outcomes.total} on \GA, all unsupported operations of \polarsgpu{} on
the pre-declared \q{cx\_q13} and \q{cx\_window}. A G8 flag disables only a session's stream
cross-check (A4); where it ran, the median relative difference of the sampled integral from the counter
was \stat{qa.crosscheck.h100.median} on \GH{} and \stat{qa.crosscheck.a100.median} on \GA.
The B1 check on \CV{} gave a time ratio VM/metal of \stat{c_vm_check.ratio_vm_over_metal} (95\%~CI
\stat{c_vm_check.ci_low}--\stat{c_vm_check.ci_high}), outside $[0.90, 1.10]$, so the check failed. As the design
prescribes, B1 therefore uses the \CV{} times where they exist (baseline SF10, SF30 and SF300, a fixed
reference per cell, not paired by round) and the metal times elsewhere (Appendix~\ref{app:readings}, item h).

% The break-even table is defined here so that it floats to the top of the page with Section 8.4.
\begin{table*}[t]
\caption{Break-even scale brackets (previous scale, $s$] of X3 under P-NEW;
``$\le$SF1'' and ``$>$SF300'' censored. Energy: lower bound on the node-level break-even, under
\NOne. Keys \texttt{break\_even.x3.<host>.<suite>.<metric>.text}.}
\label{tab:breakeven}
\footnotesize
\begin{tabular}{@{}llccccc@{}}
\toprule
Host & Suite & Cost B1 & Cost B2 & Cost B3 & Energy & Energy above idle \\
\midrule
\multirow{2}{*}{\GH} & baseline & \tstat{break_even.x3.h100.baseline.cost_b1_p_new.text} & \tstat{break_even.x3.h100.baseline.cost_b2_p_new.text} & \tstat{break_even.x3.h100.baseline.cost_b3_p_new.text} & \tstat{break_even.x3.h100.baseline.energy_total.text} & \tstat{break_even.x3.h100.baseline.energy_above_idle.text} \\
 & complex & \tstat{break_even.x3.h100.complex.cost_b1_p_new.text} & \tstat{break_even.x3.h100.complex.cost_b2_p_new.text} & \tstat{break_even.x3.h100.complex.cost_b3_p_new.text} & \tstat{break_even.x3.h100.complex.energy_total.text} & \tstat{break_even.x3.h100.complex.energy_above_idle.text} \\
\midrule
\multirow{2}{*}{\GA} & baseline & \tstat{break_even.x3.a100.baseline.cost_b1_p_new.text} & \tstat{break_even.x3.a100.baseline.cost_b2_p_new.text} & \tstat{break_even.x3.a100.baseline.cost_b3_p_new.text} & \tstat{break_even.x3.a100.baseline.energy_total.text} & \tstat{break_even.x3.a100.baseline.energy_above_idle.text} \\
 & complex & \tstat{break_even.x3.a100.complex.cost_b1_p_new.text} & \tstat{break_even.x3.a100.complex.cost_b2_p_new.text} & \tstat{break_even.x3.a100.complex.cost_b3_p_new.text} & \tstat{break_even.x3.a100.complex.energy_total.text} & \tstat{break_even.x3.a100.complex.energy_above_idle.text} \\
\bottomrule
\end{tabular}
\end{table*}

\subsection{Energy per query}
\label{sec:results-energy}
Every energy contrast here is cross-class: \scope{S2-pinned} on \CG{} over \scope{S1-counter} on the GPU host
(Section~\ref{sec:n1}; Figure~\ref{fig:energy}; Table~\ref{tab:x3} for X3 at every scale). At SF300 on \GH,
X1 is \rtext{x1.h100.baseline.sf300.energy_total} on the baseline suite and
\rtext{x1.h100.complex.sf300.energy_total} on the complex suite, and X2 is
\rtext{x2.h100.baseline.sf300.energy_total} and \rtext{x2.h100.complex.sf300.energy_total}; on \GA, X1 is
\rtext{x1.a100.baseline.sf300.energy_total} and \rtext{x1.a100.complex.sf300.energy_total}, and X2
\rtext{x2.a100.baseline.sf300.energy_total} and \rtext{x2.a100.complex.sf300.energy_total}. These and every
other estimable energy cell of X1--X3 in Table~\ref{tab:x12}, total and above idle, on both hosts, are
\ecls{x1.h100.baseline.sf300.energy_total}.

On \GH{} the tipping point of X3 at SF300 (baseline suite) is
\dstar{delta_star.x3.h100.baseline.sf300.energy_total}. On \GA{} it is
\stat{delta_star.x3.a100.baseline.sf300.energy_total.wh_per_suite_run.mean}\,Wh per suite run (95\%~CI \stat{delta_star.x3.a100.baseline.sf300.energy_total.wh_per_suite_run.ci_low} to
\stat{delta_star.x3.a100.baseline.sf300.energy_total.wh_per_suite_run.ci_high}; \stat{delta_star.x3.a100.baseline.sf300.energy_total.w_over_gpu_suite_time.mean}\,W), also non-negative. For scale, the pinned \CG{}
socket idled at \stat{qa.idle_power_w.c-gcp.polars.S2-pinned.mean}\,W and the GPU board under a CPU
engine at \stat{qa.idle_power_w.h100.polars.S1-counter.mean}\,W (H100) and
\stat{qa.idle_power_w.a100.polars.S1-counter.mean}\,W (A100). Above idle, X3 at SF300 is
\rtext{x3.h100.baseline.sf300.energy_above_idle} on \GH{} and \rtext{x3.a100.baseline.sf300.energy_above_idle}
on \GA, both \ecls{x3.a100.baseline.sf300.energy_above_idle}.

\subsection{Cost per query}
\label{sec:results-cost}

\begin{figure}[t]
\centering
% Drawn by tools/make_figures.py from stats.json (display rules in its header); placeholder sized as the figure.
\studyfig{f02_time_vs_price}{\columnwidth}{3.5in}
\caption{Factorised cost: the X3 time ratio $t_C/t_G$ against \CG{} (log scale; whiskers: 95\%
interval) per host, suite and scale, with the B1 price ratio $p_G/p_C$ under P-NEW (solid) and P-OLD
(dashed; the A100's equal). Points above a line are cheaper on the GPU host under that table where B1 uses
the metal times (not at baseline SF10, SF30 and SF300; Section~\ref{sec:results-qa}). Marker shape gives
the class of the time ratio.}
\label{fig:factorised}
\end{figure}

Under B1 at P-NEW (Table~\ref{tab:x3}), X3 at SF300 on the baseline suite is
\rtext{x3.h100.baseline.sf300.cost_b1_p_new} on \GH, \cls{x3.h100.baseline.sf300.cost_b1_p_new}, and
\rtext{x3.a100.baseline.sf300.cost_b1_p_new} on \GA, \cls{x3.a100.baseline.sf300.cost_b1_p_new}; on the complex
suite it is \rtext{x3.h100.complex.sf300.cost_b1_p_new}, \cls{x3.h100.complex.sf300.cost_b1_p_new}, and
\rtext{x3.a100.complex.sf300.cost_b1_p_new}, \cls{x3.a100.complex.sf300.cost_b1_p_new}. Against the metal host the time ratios of these baseline cells are
\rtext{x3.h100.baseline.sf300.time_vs_cgcp} and \rtext{x3.a100.baseline.sf300.time_vs_cgcp} (Figure~\ref{fig:factorised}),
against price ratios of 1.680 and 0.684 (Section~\ref{sec:cost}); the B1 cost at baseline SF300 uses the
\CV{} times instead (Section~\ref{sec:results-qa}). Within a framework (baseline
suite, SF300, B1), X1 is \rtext{x1.h100.baseline.sf300.cost_b1_p_new}, \cls{x1.h100.baseline.sf300.cost_b1_p_new},
and X2 \rtext{x2.h100.baseline.sf300.cost_b1_p_new}, \cls{x2.h100.baseline.sf300.cost_b1_p_new}, on \GH, and
\rtext{x1.a100.baseline.sf300.cost_b1_p_new}, \cls{x1.a100.baseline.sf300.cost_b1_p_new}, and
\rtext{x2.a100.baseline.sf300.cost_b1_p_new}, \cls{x2.a100.baseline.sf300.cost_b1_p_new}, on \GA. Against \CG{} as
rented (B2), X3 at SF300 is \rtext{x3.h100.baseline.sf300.cost_b2_p_new} on \GH{} and
\rtext{x3.a100.baseline.sf300.cost_b2_p_new} on \GA; against the CPU engine on the same GPU host (B3),
\rtext{x3.h100.baseline.sf300.cost_b3_p_new} and \rtext{x3.a100.baseline.sf300.cost_b3_p_new}. For P-BILLED, \GH{}
ran \stat{billing_by_host.h100.hours}\,h billed at \stat{billing_by_host.h100.billed_usd}\,USD, \GA{}
\stat{billing_by_host.a100.hours}\,h at \stat{billing_by_host.a100.billed_usd}\,USD and \CG{}
\stat{billing_by_host.c-gcp.hours}\,h at \stat{billing_by_host.c-gcp.billed_usd}\,USD (hours $\times$ price at
launch); the providers' invoices were not available, so the invoice gap is not computed (A10).

\subsection{Break-even scales}
\label{sec:results-breakeven}

Under B1 at P-NEW the X3 brackets (Table~\ref{tab:breakeven}) are
\stat{break_even.x3.h100.baseline.cost_b1_p_new.text} and \stat{break_even.x3.h100.complex.cost_b1_p_new.text} on
\GH{} (baseline, complex) and \stat{break_even.x3.a100.baseline.cost_b1_p_new.text} and
\stat{break_even.x3.a100.complex.cost_b1_p_new.text} on \GA, the same under P-OLD and P-BILLED
(Table~\ref{tab:be-all}; the A100's P-OLD price equals P-NEW). The energy brackets, lower bounds on the
node-level energy break-even under \NOne, are \stat{break_even.x3.h100.baseline.energy_total.text} on the
baseline suite and \stat{break_even.x3.h100.complex.energy_total.text} on the complex suite, on both hosts. A bracket ``$>$SF300'' means the top scale is not GPU better; ``$\le$'' the lowest scale of the suite
means every scale is; a named non-estimable scale bounds the bracket (Section~\ref{sec:classification}).

\subsection{GPU generation, CPU platform and the idle-GPU tax}
\label{sec:results-x456}
X4 (A100 over H100; descriptive, confounded with the host; above 1, the A100 needs more): at SF300 on
the baseline suite the ratios for
\polarsgpu{} are, for time, \rtext{x4.a100.polars_gpu.baseline.sf300.time}, for board energy
\rtext{x4.a100.polars_gpu.baseline.sf300.energy_total} and for cost
\rtext{x4.a100.polars_gpu.baseline.sf300.cost_p_new}, and for \sparkgpu{}
\rtext{x4.a100.spark_gpu.baseline.sf300.time}, \rtext{x4.a100.spark_gpu.baseline.sf300.energy_total} and
\rtext{x4.a100.spark_gpu.baseline.sf300.cost_p_new}. X5 (a CPU engine on the GPU host over \CG, time; above 1, slower on the GPU host): \polarscpu{}
\rtext{x5.h100.polars.baseline.sf300.time} and \sparkcpu{} \rtext{x5.h100.spark_cpu.baseline.sf300.time} on \GH,
\rtext{x5.a100.polars.baseline.sf300.time} and \rtext{x5.a100.spark_cpu.baseline.sf300.time} on \GA. X6 (the
idle-GPU tax: cost on the GPU host over B1; above 1, dearer on the GPU host): for \polarscpu{} the ratio is
\rtext{x6.h100.polars.baseline.sf300.cost_ratio_b3_over_b1.p_new} on \GH{} (price ratio
\stat{x6.h100.polars.baseline.sf300.price_ratio_b3_over_b1.p_new}) and
\rtext{x6.a100.polars.baseline.sf300.cost_ratio_b3_over_b1.p_new} on \GA{} (price ratio
\stat{x6.a100.polars.baseline.sf300.price_ratio_b3_over_b1.p_new}), and the idle GPU board used
\stat{x6.h100.polars.baseline.sf300.idle_gpu_energy_wh_per_suite_run.mean}\,Wh (\GH) and
\stat{x6.a100.polars.baseline.sf300.idle_gpu_energy_wh_per_suite_run.mean}\,Wh (\GA) per baseline suite run.

\subsection{Sensitivity}
\label{sec:results-sens}
No accepted round lies outside a complete square: the stopped launch~1s of \CV{} and \CG{} (A8)
have no accepted round, and \CG{} launch~2 was not run (A9). Per-launch estimates and the unpaired Welch estimates of every cross-host cell are in the
deposited analysis output. For X3 cost under B1 at SF300 on \GH{} the unpaired estimate is
\rtext{sensitivity_unpaired.x3.h100.baseline.sf300.cost_b1_p_new}, and the strict X3 (queries both GPU engines
ran) is \rtext{x3_strict.h100.baseline.sf300.cost_b1_p_new}.

% =====================================================================================
\section{Discussion}
\label{sec:discussion}
% One interpretive sentence per paragraph (the last of each), from the classes, intervals and brackets of
% Section 8 only; no new analysis.
\paragraph{Using the results to choose an engine and host.} The factorised form (Section~\ref{sec:cost},
Figure~\ref{fig:factorised}) lets an organization re-price every cell: a GPU host is cheaper per query
wherever the measured time ratio exceeds the organization's own price ratio, and the brackets of
Table~\ref{tab:breakeven} move with the CPU baseline (B1, B2, B3) and the price table. On this workload they moved with the baseline and not with the
price table (on the baseline suite X3 was GPU better at every scale against \CG{} as rented, and the B1
brackets are the same under all three tables; Table~\ref{tab:be-all}), and the slower A100 host was the
cheaper one per query (X4, descriptive).

\paragraph{What the energy results can and cannot say.} Every cross-class energy cell is a statement
about measured scopes. Under \NOne, only a ``CPU better'' cell supports a node-level direction; for any
other cell $\Dstar$ (Section~\ref{sec:results-energy}) is the unmeasured excess the GPU node would have to
stay below, set beside the idle power of one \CG{} socket and of the idle GPU. Here every estimable cell
was GPU better at the measured scopes, and the X3 tipping points at baseline SF300,
\stat{delta_star.x3.h100.baseline.sf300.energy_total.w_over_gpu_suite_time.mean}\,W (H100) and
\stat{delta_star.x3.a100.baseline.sf300.energy_total.w_over_gpu_suite_time.mean}\,W (A100), are several
times the \stat{qa.idle_power_w.c-gcp.polars.S2-pinned.mean}\,W idle of the pinned socket; whether the GPU
node's unmeasured host CPU and DRAM stay below them was not measured.

\paragraph{The protocol as a procedure.} Three pre-lock checks changed what could be reported: the
scope check showed the GH200 counter to cover the whole module (Appendix~\ref{app:notrun}), the
resolution probe set per-host $\tmin$ of 1--2\,s and with it the cells below resolution
(Section~\ref{sec:results-qa}), and the idle probe separated a context-open GPU from a bare one
(Section~\ref{sec:idle}). Each is a short synthetic battery per host class (Appendix~\ref{app:validation})
and should precede any per-query energy claim on rented hosts. In this campaign the resolution limit
decided where energy could be reported at all: \polarsgpu{} ran below $\tmin$ in every baseline cell up to
SF100 on both GPU hosts (analysis table t14), so the energy brackets are censored at the lowest estimable
scale, not located.

% =====================================================================================
\section{Threats to Validity and Limitations}
\label{sec:threats}

\paragraph{Virtualized counters, missing CPU energy and \NOne.} Host CPU and DRAM energy of
GPU-engine runs are unmeasured (Section~\ref{sec:scopes}); no wall-power reference exists on any host (no BMC), and RAPL covers package
and DRAM only, not power supplies, fans, storage or network. \CG's virtualization verdict was
ambiguous (no hypervisor flag; the platform reports Google), but its counters advance, agree
between powercap and \texttt{perf}, and are re-checked by G6 in every session. Its probes ran
on the same shape in \texttt{us-east4}; the campaign host is in \texttt{us-central1-a} (A5.3).
\scope{S1} and \scope{S2} are different slices of a node, and \NOne{} is untested; the tipping points show its weight without testing it. The one design element that avoided \NOne, the GH200 module-scope comparison, was not run.

\paragraph{Over-provisioned CPU host.} \CG{} has 96 cores and 1.5\,TB of DRAM, so the pinned
socket idles high against light work and its DRAM includes capacity a right-sized node lacks.
This inflates $E_C$, a bias against the CPU engines: cross-class ``CPU better'' findings are
conservative, while ``GPU better at the measured scope'' may partly reflect host size. Package and
DRAM are reported separately and above-idle beside total for this reason. B1 prices a 44-vCPU VM
for 26 threads (the vCPU surplus is not bounded) and, its time check having failed, uses the metal times
outside baseline SF10, SF30 and SF300; if the VM's gap holds there, those cells understate the CPU-side cost
and favour the CPU engines.

\paragraph{Fewer \CG{} rounds.} \CG's confirmatory estimates come from one complete square (A9: $R = 4$,
df~3). Every cross-host cell paired with \CG{} (cross-class energy, B1 and B2 cost, X5, X6) uses
rounds 1--4 only (Appendix~\ref{app:readings}, item a), so its interval rests on 3 degrees of freedom
($t = 3.182$ against 2.365 for eight rounds) and is wider for a given between-round spread. The GPU hosts'
rounds 5--8 are not in those cells; they enter the GPU hosts' within-host cells, the GPU-to-GPU contrast X4
and the labelled unpaired Welch sensitivity. The cut was made by the clock, before any result was seen.

\paragraph{H100 driver.} The two H100 squares ran different drivers, \stat{campaign.h100.driver}, so the
driver version is confounded with the instance and region of each square. It was not controlled, because
the design fixes the image and not the driver.
% Exploratory (lead's decision, 2026-11-25): computed by p3-analysis from tables/t10_cells.csv after
% unblinding and re-computed by the writer (H100 l2/l1: polars 1.116, polars_gpu 1.089, spark_cpu 1.121,
% spark_gpu 1.093; A100 1.001-1.009). Not a stats.json key.
Exploratory, computed after unblinding (2026-11-25) from the analysis output's per-cell table (t10): on
the H100 every engine, the CPU engines included, ran 9--12\% slower on launch~2 than on launch~1 (per
engine, geometric mean over cells of the launch-2/launch-1 ratio of geometric-mean round times), while the two A100 launches agreed within 1\%. A slowdown that also hits the CPU engines points
to a host- or region-level difference rather than to the GPU driver alone; the design cannot separate
the two. The X3 ratios move far less (time at baseline SF300:
\stat{sensitivity.launch:h100-l1.x3.h100.baseline.sf300.time.estimate} against
\stat{sensitivity.launch:h100-l2.x3.h100.baseline.sf300.time.estimate}), and every H100 cell paired with \CG{} rests on launch~1 alone.

\paragraph{Sensors and idle draw.} The one-sided short-run errors of the A100 and \CG{}
(Section~\ref{sec:resolution}) were measured at one power level, so a query drawing less power may
have a larger relative error (not tested).
Back-to-back window runs are warmer than interleaved trials, and thermal carryover is limited by
the Latin square and the cool-down, not removed. Idle power depends on temperature, persistence
mode, C-states and the CUDA-context state (recorded, not controlled).

\paragraph{Tenancy, prices and confounds.} Cloud hosts and their storage are shared; noise is
bounded only by rounds and two launches. Each GPU host's two launches ran concurrently on separate
instances (the design requires separation, not an order); the \CV{} rounds and \CG's first square
come from the A8 relaunches. \CG's data sat on a Hyperdisk network volume and the GPU hosts' on the
instance's virtio disk, whose backing the provider does not state; the storage types differ and are
recorded, and Spark spill I/O may differ with them. Prices are a snapshot: Lambda changed three of four
relevant prices on 2026-11-16, and spot, reserved and committed-use prices are not measured; the
factorised time and price ratios are the durable result. The \CG{} Xeon 8481C and the H100 host's 8480+ are one
family, not one SKU; the A100 host has an AMD EPYC, confounding A100 CPU-engine comparisons and
X4. One query at a time on a whole instance: real fleets share, queue and idle.

\paragraph{Not blind; scope of claims.} The prior study's timings were \camera{known to the
author}{public} before the lock, so Lambda cost ratios are partly predictable; the pre-registration fixes the choices, not the knowledge. Four
engines, one TPC-H-derived workload, five scales and one provider per host class: no claim extends
to other engines, formats, clusters, CPU vendors or query mixes, nor to GPU generations other than the
two measured.

% =====================================================================================
\section{Conclusion}
\label{sec:conclusion}
% Findings in the wording rules of Sections 4.6 and 6.4.
Energy and cost per query of CPU and GPU dataframe engines can be measured on rented hosts, but
only with the meter's scope, resolution and idle state established first and the scope carried
with every number. We locked such a protocol and its analysis plan before the campaign, with cross-class
energy claims bounded by an explicit assumption and its tipping points.
On this workload, against the right-sized CPU baseline, the cost contrast of the best CPU and GPU engines at
SF300 was \cls{x3.h100.baseline.sf300.cost_b1_p_new} on both GPU hosts
(\rtext{x3.h100.baseline.sf300.cost_b1_p_new} on the H100, \rtext{x3.a100.baseline.sf300.cost_b1_p_new} on the
A100), with break-even brackets of \stat{break_even.x3.h100.baseline.cost_b1_p_new.text} and
\stat{break_even.x3.a100.baseline.cost_b1_p_new.text} on the baseline suite and GPU better at every
complex-suite scale. Energy could be compared across classes only where runs were long enough to resolve;
there every cell was GPU better at the measured scopes, a statement about the GPU board against one
pinned CPU socket and its DRAM that supports no node-level claim.

% =====================================================================================
\section*{Artifact Availability}
\label{sec:artifacts}
% TODO(author): code release DOI for the Paper 3 harness and analysis; raw-data DOI of the
% Paper 3 campaign (both to be minted after the campaign).
The probe tools, harness, gates, analysis code, price snapshots and probe evidence are in the
study repository (\repourl); the locked design is deposited with a time stamp (\lockdoi). The
analysis code was written after the lock and frozen before any campaign data was seen (A7, freeze
commit \camera{\stat{amend.a7.freeze_commit}}{\anonmark}; synthetic recovery evidence: commit
\camera{\stat{amend.a7.recovery_commit}}{\anonmark}); the analysis output records the code's sha256
(\hashstat{_meta.analysis_code_sha256}), equal to the A7 value. Raw campaign records will be
deposited with checksums \todoplaceholder{raw-data DOI}.

\ifanon
\section*{Disclosures}
\label{sec:disclosures}
% Anonymized build only (lead's decisions, 2026-11-23). The camera-ready acknowledgments are withheld.
% (1) The AI statement appears in full: it identifies no one, ACM's policy expects AI use to be disclosed
%     in the submitted paper, and the author requires that the disclosure never omit the automation and
%     the writing assistance. It is the camera-ready text with "The author" as "The author(s)".
% (2) A credit disclosure that names no source: the CFP says "Remove acknowledgment of identifying names
%     and funding sources", and design sec. 8 discloses the sponsor relationship in every output. No
%     "no role" claim here: amendment A5 item 2 records none for the GCP credits.
The author(s) directed the study and approved every design and analysis decision. AI coding
assistants (Anthropic's Claude Code; OpenAI Codex) were used as tools to write and test code,
automate the cloud runs, and assist with writing. The author(s) reviewed all of it and take full
responsibility.

Part of the compute was paid for with cloud-provider credits; cost is computed from list prices
whatever funded the bill. The sources are named in the final version.
\fi

\ifanon\else
\begin{acks}
The author directed the study and approved every design and analysis decision. AI coding
assistants (Anthropic's Claude Code; OpenAI Codex) were used as tools to write and test code,
automate the cloud runs, and assist with writing. The author reviewed all of it and takes full
responsibility.

% Funding: design sec. 8 (Lambda, AWS, with the sponsor-role statement) and amendment A5 item 2 (GCP,
% no role statement), worded as the lead instructed on 2026-11-23.
Lambda provided \$10,000 and AWS \$6,000 in compute credits for this study (they did not cover the
runs of the two earlier studies~\cite{ejiofor2026p1,ejiofor2026p2}); neither sponsor had any role in
the design, analysis, interpretation, review or the decision to publish. GCP compute was paid by
Google Cloud ecosystem credits (2024). This study reports Lambda prices and compares Lambda GPU
hosts with GCP CPU hosts; cost uses list prices whatever funded the bill.

\end{acks}
\fi

\bibliographystyle{ACM-Reference-Format}
\bibliography{references}

% =====================================================================================
% The appendix starts on a fresh page: this flushes every pending float before it and removes the
% end-of-document overfull boxes that the appendix floats otherwise caused (tested in all four builds).
\clearpage
\appendix

\section{Full Contrast Tables}
\label{app:tables}
% Key-driven tables (generated by the writer's script from the key scheme; the values come from stats.json via
% make stats). Per-query values, totals and the coverage matrix are in the analysis output's CSV tables, deposited
% with the data.
Tables~\ref{tab:x12}--\ref{tab:be-all} complete Tables~\ref{tab:x3} and~\ref{tab:breakeven}: the X1 and X2 classes,
the 95\% intervals of X3, and the break-even brackets of X1--X3 under every baseline and price table. Per-query
values, suite totals and the coverage matrix are in the analysis output deposited with the data.

\begin{table*}[t]
\caption{X1 (\polarscpu/\polarsgpu) and X2 (\sparkcpu/\sparkgpu) per host, suite and scale; cell: estimate of $\rho$
and class as in Table~\ref{tab:x3}. Cost under P-NEW; energy at the measured scopes.}
\label{tab:x12}
\footnotesize
\setlength{\tabcolsep}{4pt}
\begin{tabular}{@{}ll*{8}{c}@{}}
\toprule
& & \multicolumn{5}{c}{baseline suite} & \multicolumn{3}{c}{complex suite} \\
\cmidrule(lr){3-7}\cmidrule(l){8-10}
Contrast & Metric & SF1 & SF10 & SF30 & SF100 & SF300 & SF30 & SF100 & SF300 \\
\midrule
\multirow{5}{*}{X1, \GH} & cost B1 & \xcell{x1.h100.baseline.sf1.cost_b1_p_new} & \xcell{x1.h100.baseline.sf10.cost_b1_p_new} & \xcell{x1.h100.baseline.sf30.cost_b1_p_new} & \xcell{x1.h100.baseline.sf100.cost_b1_p_new} & \xcell{x1.h100.baseline.sf300.cost_b1_p_new} & \xcell{x1.h100.complex.sf30.cost_b1_p_new} & \xcell{x1.h100.complex.sf100.cost_b1_p_new} & \xcell{x1.h100.complex.sf300.cost_b1_p_new} \\
 & cost B2 & \xcell{x1.h100.baseline.sf1.cost_b2_p_new} & \xcell{x1.h100.baseline.sf10.cost_b2_p_new} & \xcell{x1.h100.baseline.sf30.cost_b2_p_new} & \xcell{x1.h100.baseline.sf100.cost_b2_p_new} & \xcell{x1.h100.baseline.sf300.cost_b2_p_new} & \xcell{x1.h100.complex.sf30.cost_b2_p_new} & \xcell{x1.h100.complex.sf100.cost_b2_p_new} & \xcell{x1.h100.complex.sf300.cost_b2_p_new} \\
 & cost B3 & \xcell{x1.h100.baseline.sf1.cost_b3_p_new} & \xcell{x1.h100.baseline.sf10.cost_b3_p_new} & \xcell{x1.h100.baseline.sf30.cost_b3_p_new} & \xcell{x1.h100.baseline.sf100.cost_b3_p_new} & \xcell{x1.h100.baseline.sf300.cost_b3_p_new} & \xcell{x1.h100.complex.sf30.cost_b3_p_new} & \xcell{x1.h100.complex.sf100.cost_b3_p_new} & \xcell{x1.h100.complex.sf300.cost_b3_p_new} \\
 & energy & \xcell{x1.h100.baseline.sf1.energy_total} & \xcell{x1.h100.baseline.sf10.energy_total} & \xcell{x1.h100.baseline.sf30.energy_total} & \xcell{x1.h100.baseline.sf100.energy_total} & \xcell{x1.h100.baseline.sf300.energy_total} & \xcell{x1.h100.complex.sf30.energy_total} & \xcell{x1.h100.complex.sf100.energy_total} & \xcell{x1.h100.complex.sf300.energy_total} \\
 & above idle & \xcell{x1.h100.baseline.sf1.energy_above_idle} & \xcell{x1.h100.baseline.sf10.energy_above_idle} & \xcell{x1.h100.baseline.sf30.energy_above_idle} & \xcell{x1.h100.baseline.sf100.energy_above_idle} & \xcell{x1.h100.baseline.sf300.energy_above_idle} & \xcell{x1.h100.complex.sf30.energy_above_idle} & \xcell{x1.h100.complex.sf100.energy_above_idle} & \xcell{x1.h100.complex.sf300.energy_above_idle} \\
\midrule
\multirow{5}{*}{X1, \GA} & cost B1 & \xcell{x1.a100.baseline.sf1.cost_b1_p_new} & \xcell{x1.a100.baseline.sf10.cost_b1_p_new} & \xcell{x1.a100.baseline.sf30.cost_b1_p_new} & \xcell{x1.a100.baseline.sf100.cost_b1_p_new} & \xcell{x1.a100.baseline.sf300.cost_b1_p_new} & \xcell{x1.a100.complex.sf30.cost_b1_p_new} & \xcell{x1.a100.complex.sf100.cost_b1_p_new} & \xcell{x1.a100.complex.sf300.cost_b1_p_new} \\
 & cost B2 & \xcell{x1.a100.baseline.sf1.cost_b2_p_new} & \xcell{x1.a100.baseline.sf10.cost_b2_p_new} & \xcell{x1.a100.baseline.sf30.cost_b2_p_new} & \xcell{x1.a100.baseline.sf100.cost_b2_p_new} & \xcell{x1.a100.baseline.sf300.cost_b2_p_new} & \xcell{x1.a100.complex.sf30.cost_b2_p_new} & \xcell{x1.a100.complex.sf100.cost_b2_p_new} & \xcell{x1.a100.complex.sf300.cost_b2_p_new} \\
 & cost B3 & \xcell{x1.a100.baseline.sf1.cost_b3_p_new} & \xcell{x1.a100.baseline.sf10.cost_b3_p_new} & \xcell{x1.a100.baseline.sf30.cost_b3_p_new} & \xcell{x1.a100.baseline.sf100.cost_b3_p_new} & \xcell{x1.a100.baseline.sf300.cost_b3_p_new} & \xcell{x1.a100.complex.sf30.cost_b3_p_new} & \xcell{x1.a100.complex.sf100.cost_b3_p_new} & \xcell{x1.a100.complex.sf300.cost_b3_p_new} \\
 & energy & \xcell{x1.a100.baseline.sf1.energy_total} & \xcell{x1.a100.baseline.sf10.energy_total} & \xcell{x1.a100.baseline.sf30.energy_total} & \xcell{x1.a100.baseline.sf100.energy_total} & \xcell{x1.a100.baseline.sf300.energy_total} & \xcell{x1.a100.complex.sf30.energy_total} & \xcell{x1.a100.complex.sf100.energy_total} & \xcell{x1.a100.complex.sf300.energy_total} \\
 & above idle & \xcell{x1.a100.baseline.sf1.energy_above_idle} & \xcell{x1.a100.baseline.sf10.energy_above_idle} & \xcell{x1.a100.baseline.sf30.energy_above_idle} & \xcell{x1.a100.baseline.sf100.energy_above_idle} & \xcell{x1.a100.baseline.sf300.energy_above_idle} & \xcell{x1.a100.complex.sf30.energy_above_idle} & \xcell{x1.a100.complex.sf100.energy_above_idle} & \xcell{x1.a100.complex.sf300.energy_above_idle} \\
\midrule
\multirow{5}{*}{X2, \GH} & cost B1 & \xcell{x2.h100.baseline.sf1.cost_b1_p_new} & \xcell{x2.h100.baseline.sf10.cost_b1_p_new} & \xcell{x2.h100.baseline.sf30.cost_b1_p_new} & \xcell{x2.h100.baseline.sf100.cost_b1_p_new} & \xcell{x2.h100.baseline.sf300.cost_b1_p_new} & \xcell{x2.h100.complex.sf30.cost_b1_p_new} & \xcell{x2.h100.complex.sf100.cost_b1_p_new} & \xcell{x2.h100.complex.sf300.cost_b1_p_new} \\
 & cost B2 & \xcell{x2.h100.baseline.sf1.cost_b2_p_new} & \xcell{x2.h100.baseline.sf10.cost_b2_p_new} & \xcell{x2.h100.baseline.sf30.cost_b2_p_new} & \xcell{x2.h100.baseline.sf100.cost_b2_p_new} & \xcell{x2.h100.baseline.sf300.cost_b2_p_new} & \xcell{x2.h100.complex.sf30.cost_b2_p_new} & \xcell{x2.h100.complex.sf100.cost_b2_p_new} & \xcell{x2.h100.complex.sf300.cost_b2_p_new} \\
 & cost B3 & \xcell{x2.h100.baseline.sf1.cost_b3_p_new} & \xcell{x2.h100.baseline.sf10.cost_b3_p_new} & \xcell{x2.h100.baseline.sf30.cost_b3_p_new} & \xcell{x2.h100.baseline.sf100.cost_b3_p_new} & \xcell{x2.h100.baseline.sf300.cost_b3_p_new} & \xcell{x2.h100.complex.sf30.cost_b3_p_new} & \xcell{x2.h100.complex.sf100.cost_b3_p_new} & \xcell{x2.h100.complex.sf300.cost_b3_p_new} \\
 & energy & \xcell{x2.h100.baseline.sf1.energy_total} & \xcell{x2.h100.baseline.sf10.energy_total} & \xcell{x2.h100.baseline.sf30.energy_total} & \xcell{x2.h100.baseline.sf100.energy_total} & \xcell{x2.h100.baseline.sf300.energy_total} & \xcell{x2.h100.complex.sf30.energy_total} & \xcell{x2.h100.complex.sf100.energy_total} & \xcell{x2.h100.complex.sf300.energy_total} \\
 & above idle & \xcell{x2.h100.baseline.sf1.energy_above_idle} & \xcell{x2.h100.baseline.sf10.energy_above_idle} & \xcell{x2.h100.baseline.sf30.energy_above_idle} & \xcell{x2.h100.baseline.sf100.energy_above_idle} & \xcell{x2.h100.baseline.sf300.energy_above_idle} & \xcell{x2.h100.complex.sf30.energy_above_idle} & \xcell{x2.h100.complex.sf100.energy_above_idle} & \xcell{x2.h100.complex.sf300.energy_above_idle} \\
\midrule
\multirow{5}{*}{X2, \GA} & cost B1 & \xcell{x2.a100.baseline.sf1.cost_b1_p_new} & \xcell{x2.a100.baseline.sf10.cost_b1_p_new} & \xcell{x2.a100.baseline.sf30.cost_b1_p_new} & \xcell{x2.a100.baseline.sf100.cost_b1_p_new} & \xcell{x2.a100.baseline.sf300.cost_b1_p_new} & \xcell{x2.a100.complex.sf30.cost_b1_p_new} & \xcell{x2.a100.complex.sf100.cost_b1_p_new} & \xcell{x2.a100.complex.sf300.cost_b1_p_new} \\
 & cost B2 & \xcell{x2.a100.baseline.sf1.cost_b2_p_new} & \xcell{x2.a100.baseline.sf10.cost_b2_p_new} & \xcell{x2.a100.baseline.sf30.cost_b2_p_new} & \xcell{x2.a100.baseline.sf100.cost_b2_p_new} & \xcell{x2.a100.baseline.sf300.cost_b2_p_new} & \xcell{x2.a100.complex.sf30.cost_b2_p_new} & \xcell{x2.a100.complex.sf100.cost_b2_p_new} & \xcell{x2.a100.complex.sf300.cost_b2_p_new} \\
 & cost B3 & \xcell{x2.a100.baseline.sf1.cost_b3_p_new} & \xcell{x2.a100.baseline.sf10.cost_b3_p_new} & \xcell{x2.a100.baseline.sf30.cost_b3_p_new} & \xcell{x2.a100.baseline.sf100.cost_b3_p_new} & \xcell{x2.a100.baseline.sf300.cost_b3_p_new} & \xcell{x2.a100.complex.sf30.cost_b3_p_new} & \xcell{x2.a100.complex.sf100.cost_b3_p_new} & \xcell{x2.a100.complex.sf300.cost_b3_p_new} \\
 & energy & \xcell{x2.a100.baseline.sf1.energy_total} & \xcell{x2.a100.baseline.sf10.energy_total} & \xcell{x2.a100.baseline.sf30.energy_total} & \xcell{x2.a100.baseline.sf100.energy_total} & \xcell{x2.a100.baseline.sf300.energy_total} & \xcell{x2.a100.complex.sf30.energy_total} & \xcell{x2.a100.complex.sf100.energy_total} & \xcell{x2.a100.complex.sf300.energy_total} \\
 & above idle & \xcell{x2.a100.baseline.sf1.energy_above_idle} & \xcell{x2.a100.baseline.sf10.energy_above_idle} & \xcell{x2.a100.baseline.sf30.energy_above_idle} & \xcell{x2.a100.baseline.sf100.energy_above_idle} & \xcell{x2.a100.baseline.sf300.energy_above_idle} & \xcell{x2.a100.complex.sf30.energy_above_idle} & \xcell{x2.a100.complex.sf100.energy_above_idle} & \xcell{x2.a100.complex.sf300.energy_above_idle} \\
\bottomrule
\end{tabular}
\end{table*}

\begin{table*}[t]
\caption{X3 (best CPU engine / best GPU engine), baseline suite: estimate of $\rho$ [95\% CI]. Cost under
P-NEW; energy at the measured scopes; time against \CG{} (the time factor of B1 and B2).}
\label{tab:x3ci-base}
\scriptsize
\setlength{\tabcolsep}{4pt}
\begin{tabular}{@{}ll*{5}{c}@{}}
\toprule
Host & Metric & SF1 & SF10 & SF30 & SF100 & SF300 \\
\midrule
\multirow{6}{*}{\GH} & cost B1 & \tratioci{x3.h100.baseline.sf1.cost_b1_p_new} & \tratioci{x3.h100.baseline.sf10.cost_b1_p_new} & \tratioci{x3.h100.baseline.sf30.cost_b1_p_new} & \tratioci{x3.h100.baseline.sf100.cost_b1_p_new} & \tratioci{x3.h100.baseline.sf300.cost_b1_p_new} \\
 & cost B2 & \tratioci{x3.h100.baseline.sf1.cost_b2_p_new} & \tratioci{x3.h100.baseline.sf10.cost_b2_p_new} & \tratioci{x3.h100.baseline.sf30.cost_b2_p_new} & \tratioci{x3.h100.baseline.sf100.cost_b2_p_new} & \tratioci{x3.h100.baseline.sf300.cost_b2_p_new} \\
 & cost B3 & \tratioci{x3.h100.baseline.sf1.cost_b3_p_new} & \tratioci{x3.h100.baseline.sf10.cost_b3_p_new} & \tratioci{x3.h100.baseline.sf30.cost_b3_p_new} & \tratioci{x3.h100.baseline.sf100.cost_b3_p_new} & \tratioci{x3.h100.baseline.sf300.cost_b3_p_new} \\
 & energy & \tratioci{x3.h100.baseline.sf1.energy_total} & \tratioci{x3.h100.baseline.sf10.energy_total} & \tratioci{x3.h100.baseline.sf30.energy_total} & \tratioci{x3.h100.baseline.sf100.energy_total} & \tratioci{x3.h100.baseline.sf300.energy_total} \\
 & above idle & \tratioci{x3.h100.baseline.sf1.energy_above_idle} & \tratioci{x3.h100.baseline.sf10.energy_above_idle} & \tratioci{x3.h100.baseline.sf30.energy_above_idle} & \tratioci{x3.h100.baseline.sf100.energy_above_idle} & \tratioci{x3.h100.baseline.sf300.energy_above_idle} \\
 & time vs \CG & \tratioci{x3.h100.baseline.sf1.time_vs_cgcp} & \tratioci{x3.h100.baseline.sf10.time_vs_cgcp} & \tratioci{x3.h100.baseline.sf30.time_vs_cgcp} & \tratioci{x3.h100.baseline.sf100.time_vs_cgcp} & \tratioci{x3.h100.baseline.sf300.time_vs_cgcp} \\
\midrule
\multirow{6}{*}{\GA} & cost B1 & \tratioci{x3.a100.baseline.sf1.cost_b1_p_new} & \tratioci{x3.a100.baseline.sf10.cost_b1_p_new} & \tratioci{x3.a100.baseline.sf30.cost_b1_p_new} & \tratioci{x3.a100.baseline.sf100.cost_b1_p_new} & \tratioci{x3.a100.baseline.sf300.cost_b1_p_new} \\
 & cost B2 & \tratioci{x3.a100.baseline.sf1.cost_b2_p_new} & \tratioci{x3.a100.baseline.sf10.cost_b2_p_new} & \tratioci{x3.a100.baseline.sf30.cost_b2_p_new} & \tratioci{x3.a100.baseline.sf100.cost_b2_p_new} & \tratioci{x3.a100.baseline.sf300.cost_b2_p_new} \\
 & cost B3 & \tratioci{x3.a100.baseline.sf1.cost_b3_p_new} & \tratioci{x3.a100.baseline.sf10.cost_b3_p_new} & \tratioci{x3.a100.baseline.sf30.cost_b3_p_new} & \tratioci{x3.a100.baseline.sf100.cost_b3_p_new} & \tratioci{x3.a100.baseline.sf300.cost_b3_p_new} \\
 & energy & \tratioci{x3.a100.baseline.sf1.energy_total} & \tratioci{x3.a100.baseline.sf10.energy_total} & \tratioci{x3.a100.baseline.sf30.energy_total} & \tratioci{x3.a100.baseline.sf100.energy_total} & \tratioci{x3.a100.baseline.sf300.energy_total} \\
 & above idle & \tratioci{x3.a100.baseline.sf1.energy_above_idle} & \tratioci{x3.a100.baseline.sf10.energy_above_idle} & \tratioci{x3.a100.baseline.sf30.energy_above_idle} & \tratioci{x3.a100.baseline.sf100.energy_above_idle} & \tratioci{x3.a100.baseline.sf300.energy_above_idle} \\
 & time vs \CG & \tratioci{x3.a100.baseline.sf1.time_vs_cgcp} & \tratioci{x3.a100.baseline.sf10.time_vs_cgcp} & \tratioci{x3.a100.baseline.sf30.time_vs_cgcp} & \tratioci{x3.a100.baseline.sf100.time_vs_cgcp} & \tratioci{x3.a100.baseline.sf300.time_vs_cgcp} \\
\bottomrule
\end{tabular}
\end{table*}

\begin{table*}[t]
\caption{X3, complex suite: estimate of $\rho$ [95\% CI], as in Table~\ref{tab:x3ci-base}.}
\label{tab:x3ci-cx}
\scriptsize
\setlength{\tabcolsep}{4pt}
\begin{tabular}{@{}ll*{3}{c}@{}}
\toprule
Host & Metric & SF30 & SF100 & SF300 \\
\midrule
\multirow{6}{*}{\GH} & cost B1 & \tratioci{x3.h100.complex.sf30.cost_b1_p_new} & \tratioci{x3.h100.complex.sf100.cost_b1_p_new} & \tratioci{x3.h100.complex.sf300.cost_b1_p_new} \\
 & cost B2 & \tratioci{x3.h100.complex.sf30.cost_b2_p_new} & \tratioci{x3.h100.complex.sf100.cost_b2_p_new} & \tratioci{x3.h100.complex.sf300.cost_b2_p_new} \\
 & cost B3 & \tratioci{x3.h100.complex.sf30.cost_b3_p_new} & \tratioci{x3.h100.complex.sf100.cost_b3_p_new} & \tratioci{x3.h100.complex.sf300.cost_b3_p_new} \\
 & energy & \tratioci{x3.h100.complex.sf30.energy_total} & \tratioci{x3.h100.complex.sf100.energy_total} & \tratioci{x3.h100.complex.sf300.energy_total} \\
 & above idle & \tratioci{x3.h100.complex.sf30.energy_above_idle} & \tratioci{x3.h100.complex.sf100.energy_above_idle} & \tratioci{x3.h100.complex.sf300.energy_above_idle} \\
 & time vs \CG & \tratioci{x3.h100.complex.sf30.time_vs_cgcp} & \tratioci{x3.h100.complex.sf100.time_vs_cgcp} & \tratioci{x3.h100.complex.sf300.time_vs_cgcp} \\
\midrule
\multirow{6}{*}{\GA} & cost B1 & \tratioci{x3.a100.complex.sf30.cost_b1_p_new} & \tratioci{x3.a100.complex.sf100.cost_b1_p_new} & \tratioci{x3.a100.complex.sf300.cost_b1_p_new} \\
 & cost B2 & \tratioci{x3.a100.complex.sf30.cost_b2_p_new} & \tratioci{x3.a100.complex.sf100.cost_b2_p_new} & \tratioci{x3.a100.complex.sf300.cost_b2_p_new} \\
 & cost B3 & \tratioci{x3.a100.complex.sf30.cost_b3_p_new} & \tratioci{x3.a100.complex.sf100.cost_b3_p_new} & \tratioci{x3.a100.complex.sf300.cost_b3_p_new} \\
 & energy & \tratioci{x3.a100.complex.sf30.energy_total} & \tratioci{x3.a100.complex.sf100.energy_total} & \tratioci{x3.a100.complex.sf300.energy_total} \\
 & above idle & \tratioci{x3.a100.complex.sf30.energy_above_idle} & \tratioci{x3.a100.complex.sf100.energy_above_idle} & \tratioci{x3.a100.complex.sf300.energy_above_idle} \\
 & time vs \CG & \tratioci{x3.a100.complex.sf30.time_vs_cgcp} & \tratioci{x3.a100.complex.sf100.time_vs_cgcp} & \tratioci{x3.a100.complex.sf300.time_vs_cgcp} \\
\bottomrule
\end{tabular}
\end{table*}

\begin{table*}[t]
\caption{Break-even brackets of X1--X3 under every baseline and price table; energy: lower bound on the
node-level break-even, under \NOne.}
\label{tab:be-all}
\scriptsize
\setlength{\tabcolsep}{4pt}
\begin{tabular}{@{}lll*{7}{c}@{}}
\toprule
Contrast & Host & Suite & B1 P-NEW & B1 P-OLD & B1 P-BILLED & B2 P-NEW & B3 P-NEW & energy & above idle \\
\midrule
X1 & \GH & baseline & \tstat{break_even.x1.h100.baseline.cost_b1_p_new.text} & \tstat{break_even.x1.h100.baseline.cost_b1_p_old.text} & \tstat{break_even.x1.h100.baseline.cost_b1_p_billed.text} & \tstat{break_even.x1.h100.baseline.cost_b2_p_new.text} & \tstat{break_even.x1.h100.baseline.cost_b3_p_new.text} & \tstat{break_even.x1.h100.baseline.energy_total.text} & \tstat{break_even.x1.h100.baseline.energy_above_idle.text} \\
X1 & \GH & complex & \tstat{break_even.x1.h100.complex.cost_b1_p_new.text} & \tstat{break_even.x1.h100.complex.cost_b1_p_old.text} & \tstat{break_even.x1.h100.complex.cost_b1_p_billed.text} & \tstat{break_even.x1.h100.complex.cost_b2_p_new.text} & \tstat{break_even.x1.h100.complex.cost_b3_p_new.text} & \tstat{break_even.x1.h100.complex.energy_total.text} & \tstat{break_even.x1.h100.complex.energy_above_idle.text} \\
X1 & \GA & baseline & \tstat{break_even.x1.a100.baseline.cost_b1_p_new.text} & \tstat{break_even.x1.a100.baseline.cost_b1_p_old.text} & \tstat{break_even.x1.a100.baseline.cost_b1_p_billed.text} & \tstat{break_even.x1.a100.baseline.cost_b2_p_new.text} & \tstat{break_even.x1.a100.baseline.cost_b3_p_new.text} & \tstat{break_even.x1.a100.baseline.energy_total.text} & \tstat{break_even.x1.a100.baseline.energy_above_idle.text} \\
X1 & \GA & complex & \tstat{break_even.x1.a100.complex.cost_b1_p_new.text} & \tstat{break_even.x1.a100.complex.cost_b1_p_old.text} & \tstat{break_even.x1.a100.complex.cost_b1_p_billed.text} & \tstat{break_even.x1.a100.complex.cost_b2_p_new.text} & \tstat{break_even.x1.a100.complex.cost_b3_p_new.text} & \tstat{break_even.x1.a100.complex.energy_total.text} & \tstat{break_even.x1.a100.complex.energy_above_idle.text} \\
\midrule
X2 & \GH & baseline & \tstat{break_even.x2.h100.baseline.cost_b1_p_new.text} & \tstat{break_even.x2.h100.baseline.cost_b1_p_old.text} & \tstat{break_even.x2.h100.baseline.cost_b1_p_billed.text} & \tstat{break_even.x2.h100.baseline.cost_b2_p_new.text} & \tstat{break_even.x2.h100.baseline.cost_b3_p_new.text} & \tstat{break_even.x2.h100.baseline.energy_total.text} & \tstat{break_even.x2.h100.baseline.energy_above_idle.text} \\
X2 & \GH & complex & \tstat{break_even.x2.h100.complex.cost_b1_p_new.text} & \tstat{break_even.x2.h100.complex.cost_b1_p_old.text} & \tstat{break_even.x2.h100.complex.cost_b1_p_billed.text} & \tstat{break_even.x2.h100.complex.cost_b2_p_new.text} & \tstat{break_even.x2.h100.complex.cost_b3_p_new.text} & \tstat{break_even.x2.h100.complex.energy_total.text} & \tstat{break_even.x2.h100.complex.energy_above_idle.text} \\
X2 & \GA & baseline & \tstat{break_even.x2.a100.baseline.cost_b1_p_new.text} & \tstat{break_even.x2.a100.baseline.cost_b1_p_old.text} & \tstat{break_even.x2.a100.baseline.cost_b1_p_billed.text} & \tstat{break_even.x2.a100.baseline.cost_b2_p_new.text} & \tstat{break_even.x2.a100.baseline.cost_b3_p_new.text} & \tstat{break_even.x2.a100.baseline.energy_total.text} & \tstat{break_even.x2.a100.baseline.energy_above_idle.text} \\
X2 & \GA & complex & \tstat{break_even.x2.a100.complex.cost_b1_p_new.text} & \tstat{break_even.x2.a100.complex.cost_b1_p_old.text} & \tstat{break_even.x2.a100.complex.cost_b1_p_billed.text} & \tstat{break_even.x2.a100.complex.cost_b2_p_new.text} & \tstat{break_even.x2.a100.complex.cost_b3_p_new.text} & \tstat{break_even.x2.a100.complex.energy_total.text} & \tstat{break_even.x2.a100.complex.energy_above_idle.text} \\
\midrule
X3 & \GH & baseline & \tstat{break_even.x3.h100.baseline.cost_b1_p_new.text} & \tstat{break_even.x3.h100.baseline.cost_b1_p_old.text} & \tstat{break_even.x3.h100.baseline.cost_b1_p_billed.text} & \tstat{break_even.x3.h100.baseline.cost_b2_p_new.text} & \tstat{break_even.x3.h100.baseline.cost_b3_p_new.text} & \tstat{break_even.x3.h100.baseline.energy_total.text} & \tstat{break_even.x3.h100.baseline.energy_above_idle.text} \\
X3 & \GH & complex & \tstat{break_even.x3.h100.complex.cost_b1_p_new.text} & \tstat{break_even.x3.h100.complex.cost_b1_p_old.text} & \tstat{break_even.x3.h100.complex.cost_b1_p_billed.text} & \tstat{break_even.x3.h100.complex.cost_b2_p_new.text} & \tstat{break_even.x3.h100.complex.cost_b3_p_new.text} & \tstat{break_even.x3.h100.complex.energy_total.text} & \tstat{break_even.x3.h100.complex.energy_above_idle.text} \\
X3 & \GA & baseline & \tstat{break_even.x3.a100.baseline.cost_b1_p_new.text} & \tstat{break_even.x3.a100.baseline.cost_b1_p_old.text} & \tstat{break_even.x3.a100.baseline.cost_b1_p_billed.text} & \tstat{break_even.x3.a100.baseline.cost_b2_p_new.text} & \tstat{break_even.x3.a100.baseline.cost_b3_p_new.text} & \tstat{break_even.x3.a100.baseline.energy_total.text} & \tstat{break_even.x3.a100.baseline.energy_above_idle.text} \\
X3 & \GA & complex & \tstat{break_even.x3.a100.complex.cost_b1_p_new.text} & \tstat{break_even.x3.a100.complex.cost_b1_p_old.text} & \tstat{break_even.x3.a100.complex.cost_b1_p_billed.text} & \tstat{break_even.x3.a100.complex.cost_b2_p_new.text} & \tstat{break_even.x3.a100.complex.cost_b3_p_new.text} & \tstat{break_even.x3.a100.complex.energy_total.text} & \tstat{break_even.x3.a100.complex.energy_above_idle.text} \\
\bottomrule
\end{tabular}
\end{table*}

\section{Pre-lock Validation of the Meters}
\label{app:validation}
Table~\ref{tab:validation} gives the pre-lock probe results on which the measurement protocol of
Section~\ref{sec:energy} rests (synthetic load only).

\paragraph{RAPL path and pinning on \CG.} Before the lock, \texttt{perf stat --per-socket}
agreed with powercap within 0.05\% per socket idle and under a socket-0 burn (bare \texttt{perf
stat -a} returns the right Joules but a runtime summed over sockets, so its power column is half
the true value). Pinning was verified: socket~0 is NUMA nodes 0 and~1; the 26-vCPU pin (CPUs 0--12
and 96--108) lies on node~0; with \texttt{MPOL\_BIND} 99.96\% of an 8\,GiB allocation landed on
the bound node; under a 26-vCPU burn on socket~0, package-0 rose by 68.7\% and package-1 changed
by $-0.12$\%.

\paragraph{Resolution.} The single-run $\tmin$ never binds: a one-run window arises only
when the first run takes at least 10\,s, where the isolated error was 0.34\% (H100) and 0.22\%
(A100). On \CG{} the probe's 200\,ms series had a median spacing of 200.0\,ms, a maximum of
211\,ms and no overrun, at a sampler cost of 0.21\% of one core.

\paragraph{Idle stability and measurement overhead.} Pairs of idle
windows agreed within 2.2\% (H100), 0.23\% (A100) and 0.0065\% (\CG, whose four RAPL zones idle at
429\,W, a lower bound on wall power). The \texttt{nvidia-smi} poller changed idle draw by at most
0.72\,W, within the H100's window-to-window variation, and the RAPL sampler cost 0.39--0.48\% of
one core; both run in every session, so their draw is inside $\Pidle$ and $E$.

\begin{table*}[t]
\caption{Pre-lock validation of the meters (synthetic load; probes P1--P7 and the G8 probe; driver
570.148.08). Ratio: counter over the 20\,ms integral of the instantaneous field of the counter's
scope (10\,s idle / 20\,s load). \texttt{power.draw}: the field it equals in every sample.
Resolution: median absolute per-run error in windows mode (Section~\ref{sec:resolution}). Idle:
two 60\,s windows, bare and with a CUDA context open (ctx); \CG: four RAPL zones. G8: onset lag of
20 pulses, median / p90 / max / drift (s). Grey rows: method validation only, hosts not run.}
\label{tab:validation}
\footnotesize
\setlength{\tabcolsep}{3pt}
\begin{tabular}{@{}lllllllll@{}}
\toprule
Host & Scope tag & Ratio idle / load & \texttt{power.draw} & Resolution, \% (0.1 / 0.25 / 0.5 / 1 / 2 / 5\,s) & $\tmin$ & Idle W, bare & Idle W, ctx & G8 lag, s \\
\midrule
\GH & \scope{S1-counter} & 0.998 / 1.0002 & 1\,s trailing avg. & 9.3 / 5.5 / 1.1 / 0.78 / 0.56 / 0.57 & 1\,s & 76.3, 75.5 & 142.4, 145.9 & 0.241 / 0.313 / 0.327 / $-$0.072 \\
\GA & \scope{S1-counter} & 1.001 / 1.003 & instantaneous & 34 / 18 / 11 / 5.3 / 1.9 / 0.69 & 2\,s & 69.8, 69.6 & 75.3, 75.4 & 0.389 / 0.522 / 0.592 / +0.007 \\
\CG & \scope{S2-pinned} & -- & -- & 20.7 / 9.8 / 4.9 / 2.6 / 1.15 / 0.40 & 2\,s & 429.02, 429.05 & -- & -- \\
\midrule
\textcolor{gray}{GH200} & \textcolor{gray}{\scope{S1-module}} & \textcolor{gray}{0.9985 / 0.9982} & \textcolor{gray}{trailing avg.} & \textcolor{gray}{11.0 / 5.95 / 2.93 / 0.41 / 0.45 / 0.47} & \textcolor{gray}{1\,s} & \textcolor{gray}{162.0, 162.0} & \textcolor{gray}{209.4, 209.7} & \textcolor{gray}{0.226 / 0.297 / 0.331 / $-$0.005} \\
\textcolor{gray}{B200} & \textcolor{gray}{\scope{S1-counter}} & \textcolor{gray}{0.9986 / 1.0049} & \textcolor{gray}{trailing avg.} & \textcolor{gray}{2.77 / 5.27 / 2.93 / 0.52 / 0.25 / 0.17} & \textcolor{gray}{1\,s} & \textcolor{gray}{143.0, 143.1} & \textcolor{gray}{191.9, 192.2} & \textcolor{gray}{0.195 / 0.305 / 0.334 / $-$0.039} \\
\bottomrule
\end{tabular}
\end{table*}

\section{Hosts Validated but Not Measured}
\label{app:notrun}
GH200 and B200 were part of the locked design and were dropped by amendment A5.1 (time and
capacity). Their pre-lock probes are reported as method validation only; no study session ran on
either host.

\paragraph{GH200 (module scope).} On a Lambda \texttt{gpu\_1x\_gh200} (Neoverse-V2, 64 vCPUs, KVM
guest, driver 570.148.08) the NVML counter agreed with the integral of
\texttt{module.power.draw.instant} (0.9985 idle, 0.9982 load) and not with GPU-only
\texttt{power.draw.instant} (1.55, 1.14), and a 26-process CPU-only burn raised it from about
209\,W to 257\,W while GPU-only \texttt{power.draw} moved by 1.5\,W: the counter covers the whole
Grace--Hopper module (CPU, GPU and memory). A first resolution probe scored against the GPU-only
field never fell below a 14.7\% error at any length, a scope offset rather than a resolution
effect, and was superseded by a rerun against the module field ($\tmin = 1$\,s,
Table~\ref{tab:validation}). The enforced GPU power limit was 700\,W, although \texttt{nvidia-smi
--query-gpu power.limit} reported 900\,W (the default and maximum). This host would have allowed
the one same-scope CPU-versus-GPU energy comparison of the design (whole-module energy, no \NOne).

\paragraph{B200 (board scope).} On a Lambda \texttt{gpu\_1x\_b200\_sxm6} (Xeon Platinum 8592+, 26
vCPUs, KVM guest, no CPU energy source) the counter matched the GPU-only instantaneous field
(0.9986, 1.0049; board scope), \texttt{power.draw} equalled the trailing average in 399 of 399
samples, idle rose from 143\,W bare to 192\,W with a context open, and G8 passed. The first
resolution probe was lost to dropped SSH sessions; a rerun gave $\tmin = 1$\,s. B200 was admitted
by amendment A2.1 and then not run (A5.1).

% Appendix tables are defined here (after Appendix A) so that both land on one page with no
% overfull box in either build.
\begin{table*}[t]
\caption{Paper 3 gates (G1--G5 as in the prior study).}
\label{tab:gates}
\footnotesize
\setlength{\tabcolsep}{3pt}
\begin{tabular}{@{}R{0.12\textwidth}R{0.62\textwidth}R{0.22\textwidth}@{}}
\toprule
Gate & Rule & On failure \\
\midrule
G6 sources & counters match the probe-declared sources; none missing or non-monotone; every sampler series advancing within the session; stamps with counter readings present & session energy invalid; round fails \\
G7 coverage & trapezoid windows have $\ge$80\% of expected samples, no gap $>$1\,s & cross-check flagged \\
G8 clocks & onset lag median within $\pm$0.4\,s, drift $\le$0.2\,s (Sec.~\ref{sec:clock}) & stream values disabled \\
G9 resolution & window rule met and per-run time $\ge \tmin$ & cell below resolution \\
G10 idle & pre/post idle mean power within 10\% per scope & above-idle flagged \\
G11 thermal & throttle reasons, enforced power limit, temperature logged & flagged, not excluded \\
G12 exclusive & no foreign GPU process; nothing else scheduled on \CG{} & integrity failure \\
\bottomrule
\end{tabular}
\end{table*}

\begin{table*}[p]% own float page: three stacked appendix tables overfill a page by about 2 pt
\caption{Amendments to the locked design. ``Seen'': outcome data seen when the amendment was made.
Synthetic: smoke tests on generated data and three trivial synthetic queries, no study query.}
\label{tab:deviations}
\footnotesize
\begin{tabular}{@{}lp{0.73\textwidth}l@{}}
\toprule
Item & Change & Seen \\
\midrule
A1.1 (2026-11-22) & No price snapshot at the lock itself; one was taken the same UTC day, 1\,h\,37\,min after the tag; its Lambda prices equal P-NEW. & none \\
A1.2 & Author decision: a launch lost to a technical failure is relaunched once and restarts its square; the lost rounds are reported, never pooled, and enter no estimate. & none \\
A1.3 & SF1 functional preflight of every engine at bring-up (as in the prior study); only verdict and failure lines are read. & none \\
A1.4 & The \CG{} G12 allow-list is filled at the first \CG{} boot after review (done in A6). & none \\
A1.5, A1.8 & Hypervisor and storage type recorded; synthetic smoke data refused in a campaign. & none \\
A1.6, A1.7 & Start order (pathfinder hosts); GCP funding to be named (closed by A5.2). & none \\
A2.1 (2026-11-22) & B200 admitted after a valid resolution probe ($\tmin = 1$\,s); superseded by A5.1. & synthetic probe \\
A2.2 & The operator sees bring-up logs only through an allow-list view. & none \\
A3.1 (2026-11-22) & Smoke test 1 (H100): not ready; G8 failed, read by the lead as spurious onsets from earlier activity. & synthetic \\
A3.2 & G8 onsets count only after a quiet GPU; the runner waits for 2\,s of quiet before each pulse; thresholds unchanged. & synthetic \\
A4.1--A4.3 (2026-11-23) & Smoke test 2: host fault (CUDA error 802); a CUDA-readiness gate was added, and a host failing it before round 1 is replaced without using the A1.2 relaunch (lead's reading). Smoke test 3: harness works; G8 flagged 2 of 4 sessions. & synthetic \\
A4.4 & Stated before campaign data: many sessions may carry a G8 flag; flags are reported and never stop a round. & none \\
A5.1 (2026-11-23) & Author decision: GH200 and B200 not run (time and capacity only); RQ4 covers A100 and H100; RQ6 (GH200 module scope) not run. & none \\
A5.2 & \camera{GCP compute (\CG, \CV) paid by Google Cloud credits}{GCP compute funding recorded (source withheld for review)}; cost still uses P-NEW. & none \\
A5.3 & No H100 capacity at the start: A100 and \CG{} (\texttt{us-central1-a}, a pre-registered fallback zone, same shape and price) started first. & none \\
% A6 checked against the committed text (commit 9a8b07b) on 2026-11-23.
A6 (2026-11-23) & \CG{} G12 allow-list after review: the 25 scheduled items found at bring-up (18 systemd timers, 7 \texttt{/etc/cron.*} files) are all masked or inert and none may run; G12 fails on any other scheduled item; same value for both \CG{} launches. & none \\
% A7: date, commits and sha256 are in record.tex (lead, 2026-11-23; amendment A7 records the
% commit, the sha256 and the synthetic-recovery result before any campaign aggregate is computed or
% seen). Checked against the committed A7 (amendment commit 5382101, freeze commit ea67f12) on 2026-11-23.
A7 (\stat{amend.a7.date}) & Contrary to design \S5.4, no analysis code existed at the lock. It was written after the lock by a team member without access to campaign data, validated only on synthetic campaigns with injected known effects, independently reviewed, and frozen before any campaign data was seen (freeze commit \camera{\stat{amend.a7.freeze_commit}}{\anonmark}; a first freeze the same day was superseded by one change to a diagnostic); code sha256 \hashstat{amend.a7.sha256}. Synthetic recovery passes (A7.3); the readings of the design it applies are in Appendix~\ref{app:readings}. & none \\
% A8: values in record.tex (lead, 2026-11-24); checked against the committed A8 (commit 50eea45).
A8 (\stat{amend.a8.date}) & Defect in the round gate of the CPU-only hosts: G1 required an idle-GPU record that is written only for jobs with a GPU engine, so every \CV{} and \CG{} session failed and both launch~1s stopped by rule. Fixed (commit \camera{\stat{amend.a8.fix_commit}}{\anonmark}); each launch relaunched once under A1.2, \CV{} first; the stopped rounds are reported, never pooled; the frozen analysis code and its hash are unaffected. Seen: only failing check ids and the QA lines. & none \\
% A9: values in record.tex (lead, 2026-11-24); checked against the committed A9 (commit 2a1727c).
A9 (\stat{amend.a9.date}) & Author decision, for time alone: \CG{} launch~2 (rounds 5--8) is not run. Under the design's square rule, \CG's confirmatory estimates come from one complete square (the relaunch, rounds 1--4), $R = 4$ (df~3), and every cell pairing \CG{} with another host uses rounds 1--4 only; the other hosts are unchanged. If that square fails for a technical reason, launch~2 is run after all. Seen: progress only (QA lines, status views). & none \\
% A10, A11: values in record.tex (lead, 2026-11-24); checked against the committed text (cb22dd8, eb097ba).
A10 (\stat{amend.a10.date}) & The analysis command and its flags are fixed before unblinding: one run of the frozen code once every launch is complete or stopped, with \CG{} launch~2 declared ended (A9); a design-hash relaxation only for the A6 allow-list value, with its expected digests; no status override; any other refusal means stop and amend. & none \\
A11 (\stat{amend.a11.date}) & Not computed, because the frozen code implements the confirmatory analyses and the stated sensitivities only: session energy and cost and cold-job cost (secondary), EX1--EX3 (exploratory, including the MODEL host-CPU estimate) and X7 (A5). None is confirmatory; the raw data are deposited. & none \\
\bottomrule
\end{tabular}
\end{table*}

\section{Readings of the Design Applied by the Analysis Code}
\label{app:readings}
Amendment A7 (item 4) fixes, before any aggregate was seen, how the frozen analysis code reads
parts of the design that its text leaves open. Each is reported here.
\begin{enumerate}[label=(\alph*)]
\item Cross-host contrasts (cross-class energy X1--X3, B1 and B2 cost, time against \CG, X5, X6)
  pair the two hosts' rounds by round number, a convention because they are independent runs; where
  the hosts share no round number the cell is not estimable, and an unpaired Welch estimate with
  each side's own rounds is given beside every cross-host cell as a labelled sensitivity.
\item X3 takes the best value among the engines that ran the query; if an engine that ran has no
  value for the metric (below resolution, an invalid counter), that query and round is unavailable
  for the metric, with no other engine standing in. The strict X3 also needs both GPU engines to
  have run; a lenient variant is a labelled diagnostic that carries no claim.
\item Break-even: a scale that was not measured is not ``not GPU better''. If the walk down from
  the top ends at a scale with no estimate, the bracket is censored (``$\le$SF$s$'', the
  unestimable scale named); if the top scale has no estimate it is ``not determinable''; and
  ``$>$SF300'' only when the top scale has a class that is not GPU better.
\item G9 drops a cell's energy only; its time and cost stay.
\item Class precedence follows Section~\ref{sec:classification}: an interval inside the
  $\pm$10\% band but wholly above 1 is GPU better, flagged as within the margin.
\item Suite-level geometric means use, per round, the queries both sides have in that round;
  suite totals, $\Dstar$ and X6 use the queries every round has; per-query ratios with a
  non-positive value (above-idle energy can be negative) are excluded and counted.
\item A complete square is all its rounds accepted within one instance; a relaunch after a loss
  (A1.2) is never pooled with the lost instance; launches of one host on different instance types
  are refused.
\item If the \CV{} check fails, B1 uses the \CV's geometric-mean time per cell as a fixed
  reference, not paired by round, and the P-BILLED B1 price is the \CV{} launch's billed price
  (P-NEW's if there is none).
\item Blinding guard: the full analysis refuses unless every planned launch (A100 and H100
  launches 1 and 2, \CG{} launches 1 and 2, \CV{} launch 1) is complete or stopped, or declared
  ended by label; until then only a QA-only mode runs, reading the gate files the QA view allows.
  The code drops X7 (A5), computes no model energy value and raises an error on an untagged energy
  value.
\end{enumerate}

\section{Integrity and Telemetry Gates}
\label{app:gates}
Table~\ref{tab:gates} gives the rule and the consequence of each Paper 3 gate (Section~\ref{sec:gates}).

\section{Pre-registration Record: Amendments A1--A7 and What Each Saw}
\label{app:amendments}
Table~\ref{tab:deviations} lists every amendment to the locked design with the outcome data seen
when it was made. It is the record on which the confirmatory status of the claims in the body
rests (Section~\ref{sec:deviations}).

% Statistics keys: paper3/README.md, section "Statistics keys" (make readme-keys).

\end{document}
